%!TEX root = ../paper.tex
%
% Revised statement of the WellFormed invariants.
%
% Purpose: a single precise specification of WellFormed, suitable as input to
%   (a) the Iris/Rocq mechanization, and
%   (b) the Clang frontend type checker.
%
% This file supersedes the invariant figures in Section \ref{sec:memaxioms} of
% chapters/RCUTypesSoundnessLemmasESOP.tex, and is now part of the paper: it is
% the first appendix section, ahead of the soundness lemmas, whose proof cases
% cite the revised invariants. Section \ref{sec:inv-impact} records which cases
% each change touched and the state of each.
%
\section{\textsf{WellFormed}, Revised}
\label{sec:inv-revised}

\newcommand{\obs}{\textsf{obs}}
\newcommand{\itr}{\textsf{iterator}}
\newcommand{\unlk}{\textsf{unlinked}}
\newcommand{\frsh}{\textsf{fresh}}
\newcommand{\frbl}{\textsf{freeable}}
\newcommand{\rt}{\textsf{root}}
\newcommand{\undf}{\textsf{undef}}
\newcommand{\Loc}{\textsf{Loc}}
\newcommand{\TID}{\textsf{TID}}
\newcommand{\FType}{\textsf{FType}}
\newcommand{\RCU}{\textsf{RCU}}
\newcommand{\defn}{\;\triangleq\;}

\subsection{Preliminaries}
\label{sec:inv-prelim}

Four notational problems recur across the original figures and are fixed once
here rather than nineteen times below.

\paragraph{(C0) Every observation carries the observing thread.}
The original grammar tags only \itr:
\[\obs ::= \itr\;\mathit{tid} \mid \unlk \mid \frsh \mid \frbl \mid \rt\]
This makes three invariants unstatable or ill-typed.  \textbf{RITR} is
\emph{about} which threads may make which observations, and cannot be phrased
when the observations in question are anonymous; \textbf{WULK} works around this
by writing $\textsf{undef}\notin O(o)$, but \textsf{undef} is not in the grammar
at all (undefinedness lives in $U$, not $O$); \textbf{WF} recovers the observer
only indirectly, through a variable that happens to point at the location.  We
therefore tag uniformly:
\[\obs ::= \itr\;\mathit{tid} \mid \unlk\;\mathit{tid} \mid \frsh\;\mathit{tid} \mid \frbl\;\mathit{tid} \mid \rt\]
\rt{} stays anonymous: it is a property of the structure, not of an observer.

This is the one deliberate \emph{design} change in this document; everything
else is a repair.  It is cheap: the type denotations in Figure
\ref{fig:denotingtypeenviroment} already conjoin $l = \mathit{tid}$ wherever an
untagged observation appears, so adding the tag introduces no new content
there --- it only makes explicit what those conjuncts already fix.  For the Iris
encoding the cost is zero, since $O$ is a $\textsf{gmap}\;\Loc\;(\textsf{gset}\;\obs)$
either way.  If this change is rejected, \textbf{RITR} and \textbf{WULK} must
instead be restated as invariants over the transition relation rather than over
a single state, which is strictly more painful to mechanize.

\paragraph{(C1) $\textsf{OID}$.}
Used in \textbf{OW} but never defined anywhere in the paper.  Read as \Loc:
the invariant is about fields holding \emph{object references}, so the guard
distinguishes pointer-valued fields from scalar ones.  We write $v\in\Loc$.

\paragraph{(C2) $h^{*}$ is partial.}
$h^{*} : (\Loc\times\textsf{Path}) \rightharpoonup \textsf{Val}$ is declared
partial in Figure \ref{fig:denotingtypeenviroment}, but several invariants
quantify over all $\rho$ as though it were total, which makes them false for
paths that run off the structure.  We write $h^{*}(\sigma.rt,\rho) = \lfloor o
\rfloor$ for ``$h^{*}$ is defined at $\rho$ and yields $o$''.

\paragraph{(C3) Free-list lookup.}
The original writes $\sigma.F([o\mapsto T])$ as a \emph{predicate} meaning
``$o\in dom(F)$ and $F(o)=T$'', and in two places binds it under
$\forall_{T'\subseteq T}$, which reads as asserting that $F$ maps $o$ to every
subset of $T$.  Worse, the bound name $T$ collides with the thread-set component
$T$ of the logical state $(\sigma,O,U,T,F)$ --- a capture bug in \textbf{FLR}
and \textbf{RINFL}.  We write $\sigma.F(o) = T_{r}$, reserving $T$ for the
state component.

\subsection{The invariants}
\label{sec:inv-list}

$\textsf{WellFormed}(\sigma,O,U,T,F)$ is the conjunction of the nineteen
predicates below.  Each carries a severity: \textsc{critical} means the original
is vacuous, unprovable, unstatable, or too weak to support the use the proofs
make of it, so a proof case discharged against it establishes nothing;
\textsc{major} means ill-formed but with recoverable intent; \textsc{minor}
means a typo, an unused binder, or a caption that contradicts its formula.

All nineteen are mechanized in \texttt{rocq/WellFormed.v}, which also carries
machine-checked proofs of the \textsc{critical} claims and a satisfiability
witness (\S\ref{sec:inv-impact}).  The development is axiom-free, and this is
checked rather than asserted: a successful build proves less than it looks,
since \texttt{Admitted} compiles and \texttt{Print Assumptions} reports an axiom
on standard output without failing anything.  \texttt{make check} rebuilds from
scratch, rejects \texttt{Admitted} and axiom declarations outright, and counts
the \texttt{Print Assumptions} reports: all 328 say \emph{closed under the
global context}, and none reports an axiom.

That target had the bug it exists to catch, and we report it because the
argument for having it is exactly what the bug illustrates.  A
\texttt{Print Assumptions} inside a \texttt{Section} prints nothing at all.
Five of ours were inside one --- both read-side triples and the whole of the
binary search tree chain --- so five results the development claimed were
checked were not being checked.  The results were sound, which we confirmed by
running the queries outside their sections; the checking was not.  They are
moved out, and \texttt{make check} now fails if an indented one appears.  The
count rose from 221 to 226 for that reason and no other: nothing new was proved,
five things merely began to be verified that had been asserted.

\paragraph{1.\ OW --- No Sharing.\hfill\textsc{minor}}
\[
\begin{array}{l}
\textbf{OW} \defn \forall o,o',f,f',v \ldotp
  \sigma.h(o,f) = v \land \sigma.h(o',f') = v \land v\in\Loc \\
\qquad {} \land \FType(f) = \RCU \land \FType(f') = \RCU \implies \\
\qquad\quad (o = o' \land f = f')
  \;\lor\; \textsf{Detached}(O,o) \;\lor\; \textsf{Detached}(O,o')
\end{array}
\]
where $\textsf{Detached}(O,o) \defn \exists \mathit{tid} \ldotp
\{\unlk\,\mathit{tid}, \frbl\,\mathit{tid}, \frsh\,\mathit{tid}\} \cap O(o) \neq \emptyset$.

\emph{Changed:} the caption (``none of the heap nodes can be observed as
undefined'') describes a different invariant than the formula, which is
in-degree-at-most-one for live nodes.  Renamed accordingly --- ``Ownership'' is
misleading.  $\FType(f') = \RCU$ added: the original constrains only $f$, which
would let a non-\RCU{} field alias a live node.  $\textsf{OID}$ resolved per (C1).

\paragraph{2.\ RWOW --- Reader/Writer Observation Coexistence.\hfill\textsc{major}}
\[
\begin{array}{l}
\textbf{RWOW} \defn \forall x,\mathit{tid},o \ldotp
  \sigma.s(x,\mathit{tid}) = o \land (x,\mathit{tid}) \notin U \implies \\
\qquad \itr\,\mathit{tid} \in O(o) \;\lor\; \big(\sigma.l = \mathit{tid} \land
  ( \unlk\,\mathit{tid} \in O(o) \lor \frbl\,\mathit{tid} \in O(o) \lor \frsh\,\mathit{tid} \in O(o) )\big)
\end{array}
\]
\emph{Changed:} parenthesisation only.  The original nests
$(\sigma.l{=}\mathit{tid} \land (\unlk \in O(o)) \lor (\sigma.l{=}\mathit{tid} \land \frbl \in O(o)))$
with unbalanced grouping, so the intended four-way disjunction does not parse as
written.

\paragraph{3.\ AWRT --- Alias With Root.\hfill\textsc{minor}}
\[
\textbf{AWRT} \defn \forall y,\mathit{tid} \ldotp
  \sigma.s(y,\mathit{tid}) = \sigma.rt \land (y,\mathit{tid}) \notin U
  \implies \itr\,\mathit{tid} \in O(\sigma.rt)
\]
\emph{Changed:} $h^{*}(\sigma.rt,\epsilon)$ replaced by $\sigma.rt$ (they are
equal by definition, and the indirection obscures the invariant).  $s$
qualified as $\sigma.s$, consistently with the rest.  $(y,\mathit{tid}) \notin
U$ added --- without it a dangling variable that happens to hold $rt$ forces an
observation into $O$.

\emph{Note for mechanization:} this makes $\sigma.rt$ carry \emph{both}
$\rt$ and $\itr\,\mathit{tid}$.  That is intended --- \textsc{T-Root} hands the
reader an $\textsf{rcuItr}\,\epsilon$ --- but $\llbracket x : \textsf{rcuRoot}
\rrbracket$ asserts $\rt \in O(rt)$, so the two must be shown to coexist rather
than to conflict.  State it as an explicit lemma.

\paragraph{4.\ IFL --- Iterators in Free List.\hfill\textsc{major}}
\[
\textbf{IFL} \defn \forall \mathit{tid},o,T_{r} \ldotp
  \itr\,\mathit{tid} \in O(o) \land \sigma.F(o) = T_{r} \implies \mathit{tid} \in T_{r}
\]
\emph{Changed:} per (C3).  Caption also overshoots: the formula concludes
$\mathit{tid} \in T_{r}$, not $\mathit{tid} \in \sigma.B$.  The latter follows,
but only via \textbf{RINFL}; the caption should say so rather than assert it
directly.

\paragraph{5.\ ULKR --- Unlinked Reachability.\hfill\textsc{critical}}
\[
\textbf{ULKR} \defn \forall o,o',f',\mathit{tid} \ldotp
  \big(\unlk\,\mathit{tid} \in O(o) \lor \frbl\,\mathit{tid} \in O(o)\big)
  \land \sigma.h(o',f') = o
  \implies \unlk\,\mathit{tid} \in O(o') \lor \frbl\,\mathit{tid} \in O(o')
\]
\emph{Changed:} the hypothesis is closed under the disjunction.  The published
version takes only $\unlk$ on the left but concludes $\unlk \lor \frbl$, so its
conclusion is weaker than its hypothesis and \textbf{it does not chain}: the
induction on path length stalls at the first $\frbl$ predecessor, where no
invariant applies.

\emph{Why this is critical:} the caption, and the prose in Section
\ref{sec:lemmas}, both state the \emph{reachability} consequence --- ``unlinked
nodes are reachable only from other unlinked nodes, and never from the root'' ---
and that is the property carrying the memory-safety argument, since it is what
stops a newly arrived reader from reaching an unlinked node.  The one-step
formula does not yield it.  With the hypothesis closed as above the induction
goes through and the reachability form is a genuine corollary.

This one was not found by reading.  It surfaced when the Rocq proof of the
transitive form would not close; the counterexample is a chain $2 \to 1 \to 0$
with $0$ unlinked, $1$ freeable --- which the published invariant explicitly
permits, since it offers exactly that disjunct --- and $2$ an ordinary iterator
reaching an unlinked node.  Both the repair and the counterexample are machine
checked: see \texttt{ULKR\_closed\_reaches} and \texttt{ULKR\_does\_not\_chain}
in \texttt{rocq/WellFormed.v}.

\paragraph{6.\ FLR --- Free List Reachability.\hfill\textsc{major}}
\[
\textbf{FLR} \defn \forall o,o',f',T_{r} \ldotp
  \sigma.F(o) = T_{r} \land \sigma.h(o',f') = o
  \implies \exists T_{r}' \ldotp \sigma.F(o') = T_{r}' \land T_{r} \bowtie T_{r}'
\]
\emph{Changed:} per (C3), including the $T$ capture.  $\bowtie$ is $\supseteq$:
$T_{r}' \subseteq T_{r}$, the predecessor's reader set contained in the
successor's, which is the direction the report already writes.  It was briefly
left open; the argument that it is in fact forced is in
Section \ref{sec:inv-open}.

\paragraph{7.\ WULK --- Writer Unlink.\hfill\textsc{major}}
\[
\textbf{WULK} \defn \forall o \ldotp \itr\,\sigma.l \in O(o) \implies
  \forall \mathit{tid} \ldotp \unlk\,\mathit{tid} \notin O(o) \land \frbl\,\mathit{tid} \notin O(o)
\]
\emph{Changed:} the caption (``the writer thread cannot observe a heap location
as \unlk'') states something the formula does not, and something that is plainly
false --- the writer observes locations as \unlk{} constantly; that is the
entire point of the type. The formula's actual content, which is the useful one,
is that \itr{} and \unlk/\frbl{} observations of the same location are mutually
exclusive.  Caption corrected. $\textsf{undef} \notin O(o)$ dropped: not in the
grammar, and subsumed by $U$.

\paragraph{8.\ FR --- Fresh Unreachable.\hfill\textsc{critical}}
\[
\begin{array}{l}
\textbf{FR} \defn \forall \mathit{tid},x,o \ldotp
  \sigma.s(x,\mathit{tid}) = o \land \frsh\,\mathit{tid} \in O(o) \implies \\
\qquad (\forall o',f' \ldotp \sigma.h(o',f') \neq o)
\;\land\;
  (\forall y,\mathit{tid}' \ldotp (y,\mathit{tid}') \neq (x,\mathit{tid}) \implies \sigma.s(y,\mathit{tid}') \neq o)
\end{array}
\]
\emph{Changed:} the original is a \emph{disjunction}
$(h(o',f') \neq o) \lor (s(y,\mathit{tid}) \neq o \lor (\mathit{tid}' \neq
\mathit{tid} \implies s(y,\mathit{tid}') \neq o))$ whose second disjunct
contradicts its own hypothesis: instantiating $y := x$ gives
$\sigma.s(x,\mathit{tid}) \neq o$, which the antecedent has just asserted is
false.  The two clauses are independent obligations --- no heap edge into a
fresh node, and no second variable aliasing it --- so this must be a
conjunction, with the alias clause excluding the witness $x$ itself.

\emph{Why this is critical:} \textbf{FR} is the sole justification for the
freshness/uniqueness step in the \textsc{Replace} and \textsc{Insert} lemmas
(both cite \texttt{ahu16f.FR} / \texttt{ahu16flk.FR} for ``there exists no field
alias/path alias to the freshly allocated $o_n$'').  In its published form it
does not support that step.

\paragraph{9.\ WF --- Fresh Only by Writer.\hfill\textsc{minor}}
\[
\textbf{WF} \defn \forall \mathit{tid},x,o \ldotp
  \sigma.s(x,\mathit{tid}) = o \land \frsh\,\mathit{tid} \in O(o) \implies \mathit{tid} = \sigma.l
\]
\emph{Note:} under (C0) this is close to definitional.  Keep it as a named
invariant for readability, but expect it to discharge in one step.

\paragraph{10.\ FNR --- Fresh Carries No Other Observation.\hfill\textsc{major}}
\[
\textbf{FNR} \defn \forall o,\mathit{tid},\mathit{tid}' \ldotp
  \frsh\,\mathit{tid} \in O(o) \implies
  \itr\,\mathit{tid}' \notin O(o) \land \unlk\,\mathit{tid}' \notin O(o)
  \land \frbl\,\mathit{tid}' \notin O(o)
\]
\emph{Changed:} two things.  (i) The original binds $x$ and never uses it.
(ii) The third conjunct is new, and it is not bookkeeping: without it the
invariant set cannot prove that \textsc{T-Replace} and \textsc{T-Insert}
preserve \textbf{ULKR}.  Both write $o_p.f := o_n$ with $o_n$ typed
\textsf{rcuFresh}; the write creates the edge $o_p \to o_n$, and \textbf{ULKR}
asks that if $o_n$ is unlinked or freeable then $o_p$ is too.  Since $o_p$ is
the writer's iterator, \textbf{WULK} denies it both observations, so the
obligation can be discharged only by denying both to $o_n$.  The published
\textbf{FNR} denies the first and says nothing about the second: \frbl{}
appears in \textbf{Detached}, \textbf{RWOW}, \textbf{ULKR}, \textbf{WULK} and
\textbf{RITR}, and in none of them as a consequent a \frsh{} node could
trigger.  So the proof does not close, and what is missing is an invariant
rather than a rule premise --- which is why the repair belongs here.

That the reachable states have no such node is beside the point, and is in fact
why the omission survived: a node acquires \frbl{} only at \textsc{SyncStop},
from \unlk{}, which the second conjunct already denies to a fresh node, so the
property does hold inductively over whole executions.  But an invariant set
earns its keep by making each action lemma provable from the invariants
\emph{alone}, and this one is not.  Adding the conjunct restores that at no
cost: \texttt{sync\_stop\_preserves\_FNR} discharges the new conjunct from the
old one, by exactly that argument.

\paragraph{11.\ FPI --- Fresh Points to Iterator.\hfill\textsc{critical}}
\[
\textbf{FPI} \defn \forall o,f,o',\mathit{tid} \ldotp
  \frsh\,\mathit{tid} \in O(o) \land \sigma.h(o,f) = o' \land \FType(f) = \RCU
  \implies \itr\,\sigma.l \in O(o')
\]
\emph{Changed:} two independent defects.  (i) The original writes
$(\frsh \in O(o) \land \exists_{f,o'} \ldotp h(o,f) = o') \implies (\forall
\mathit{tid} \ldotp \itr\,\mathit{tid} \in O(o'))$, in which $o'$ is bound by an
existential \emph{inside the antecedent} and is therefore out of scope in the
consequent.  Both must be universally bound at the top.  (ii) $\forall
\mathit{tid} \ldotp \itr\,\mathit{tid} \in O(o')$ asserts that \emph{every}
thread in the system observes $o'$ as an iterator, which is far stronger than
anything the type system establishes and is not provable.  The matching
denotation $\llbracket \textsf{rcuFresh}\,\mathcal{N} \rrbracket_{tid}$ requires
only $\itr\,\mathit{tid} \in O(s(x_i,\mathit{tid}))$ for the writer, so the
consequent is weakened to $\sigma.l$.  $\FType(f) = \RCU$ added to exclude
scalar fields.

\paragraph{12.\ WNR --- Writer Is Not a Reader.\hfill\textsc{minor}}
\[\textbf{WNR} \defn \sigma.l \notin \sigma.R\]
\emph{Unchanged.}  Well-formed as written, since $l : \TID \uplus
\{\textsf{unlocked}\}$ and $\textsf{unlocked} \notin \TID$.

\paragraph{13.\ RITR --- Readers Make Only Iterator Observations.\hfill\textsc{critical}}
\[
\textbf{RITR} \defn \forall o,\mathit{tid} \ldotp \mathit{tid} \in \sigma.R \implies
  \unlk\,\mathit{tid} \notin O(o) \land \frbl\,\mathit{tid} \notin O(o) \land \frsh\,\mathit{tid} \notin O(o)
\]
\emph{Changed:} the original reads
$\forall_{\mathit{tid} \in \sigma.R,\,o} \ldotp \itr\,\mathit{tid} \in O(o)$,
which asserts that every reader observes \emph{every} heap location as an
iterator.  That is not merely stronger than intended, it is false in any state
with an unlinked node, and it makes \textbf{WellFormed} unsatisfiable at exactly
the moments the proof needs it.  The caption gives the intent --- readers make
no observation other than \itr{} --- which is what is stated above, and which is
only expressible under (C0).

\paragraph{14.\ RINFL --- Readers in Free List Are Bounding.\hfill\textsc{major}}
\[\textbf{RINFL} \defn \forall o,T_{r} \ldotp \sigma.F(o) = T_{r} \implies T_{r} \subseteq \sigma.B\]
\emph{Changed:} per (C3), including the $T$ capture against the state's thread
set.

\paragraph{15.\ HD --- Heap Domain Closure.\hfill\textsc{critical}}
\[\textbf{HD} \defn \forall o,f,o' \ldotp \sigma.h(o,f) = o' \land o' \in \Loc
   \land \lnot\textsf{Detached}(O,o) \implies o' \in dom(\sigma.h)\]
\emph{Changed:} the original binds $f'$ and then uses $f$; $o' \in \Loc$ added to
exclude scalar-valued fields, which are otherwise required to be in the domain
of the heap; and the source is restricted to nodes that are not detached.

\emph{Why this is critical:} without the last change the invariant is false, and
\texttt{Free} is what falsifies it.  \textsc{T-UnlinkH} leaves the unlinked node
still pointing at its old child through $f_2$ --- the rule nulls its \emph{other}
fields, not that one --- and the \texttt{Free} semantics sets only the freed
node's own fields to undefined, not its predecessors'.  Nothing orders the
reclamation of two unlinked nodes, so a well-typed critical section may free a
node while another node, itself unlinked and awaiting reclamation, still points
at it.

The pointer dangles harmlessly: the node holding it is detached and unreachable,
so nothing can follow it.  That is exactly the exemption \textbf{OW} already
carries --- its conclusion is excused when either endpoint is
$\textsf{Detached}$ --- and that \textbf{OW} has it while \textbf{HD} does not
reads as an oversight rather than a decision.  Both the counterexample and the
repaired case are machine checked: see \texttt{HD\_not\_preserved\_by\_free} and
\texttt{free\_preserves\_HD}.

\emph{Note:} this is also the invariant whose proof case is \textbf{empty} in
four separate lemmas (\textsc{Replace}, \textsc{ReadHeap},
\textsc{WriteFreshField}, \textsc{Alloc}) and dismissed as trivial in a fifth
(\textsc{Free}) --- the one where it does not hold.

\paragraph{16.\ UNQRT --- Unique Root.\hfill\textsc{major}}
Split into two, because the original conflates them:
\[
\textbf{UNQRT-a} \defn \forall o,f \ldotp \sigma.h(o,f) \neq \sigma.rt
\]
\[
\textbf{UNQRT-b} \defn \forall \rho,o \ldotp h^{*}(\sigma.rt,\rho) = \lfloor o \rfloor
  \implies \itr\,\sigma.l \in O(o) \lor \rt \in O(o)
\]
\emph{Changed:} the original has a free $\mathit{tid}$; asserts
$\itr\,\mathit{tid} \in O(h^{*}(\sigma.rt,\rho))$ for all $\rho \neq \epsilon$
without regard to whether $h^{*}$ is defined there (see C2); and offers no
disjunct for $\rho = \epsilon$, where the reached node is the root and carries
\rt{} rather than \itr.  \textbf{UNQRT-a} is the no-incoming-edges half, stated
directly rather than through $h^{*}$.

\paragraph{17.\ WUNLK --- Unlinking is the Writer's.\hfill\textsc{major}}
\[
\textbf{WUNLK} \defn \forall o,\mathit{tid},\mathit{tid}_w \ldotp
  \sigma.l = \mathit{tid}_w \land
  (\unlk\,\mathit{tid} \in O(o) \lor \frbl\,\mathit{tid} \in O(o))
  \implies \mathit{tid} = \mathit{tid}_w
\]
\emph{New.}  The counterpart of \textbf{WF} at the other end of a node's life:
\textbf{WF} says an allocation is the lock holder's, and nothing said the same
of an unlinking.  It is needed, and needed locally.  The \textbf{ULKR} case of
\textsc{T-UnlinkH} and \textsc{T-Replace} ends at a detached predecessor of the
node being unlinked --- \textbf{OW} supplies the detachment, and the repaired
premise of those rules excludes a \textsf{rcuFresh} one --- so the predecessor
is unlinked or freeable by \emph{some} thread.  \textbf{ULKR}'s conclusion names
the thread its hypothesis names, so the case does not close without knowing
which.

As with \textbf{FNR}, the reachable states satisfy it anyway: \unlk{} is
produced only by the two writer rules and carried to \frbl{} by
\textsc{SyncStop}, so the thread is always the lock holder.  And as with
\textbf{FNR} that is a property of whole executions rather than an invariant, so
the action lemma does not close over it.  Guarded by the lock, so it says
nothing about a state with no writer --- the weakest form that does the job.

\paragraph{18.\ WITR --- Iterators are the Writer's or a Reader's.\hfill\textsc{critical}}
\[
\textbf{WITR} \defn \forall o,\mathit{tid} \ldotp
  \itr\,\mathit{tid} \in O(o) \implies
  \sigma.l = \mathit{tid} \lor \mathit{tid} \in \sigma.R
\]
\emph{New, and the one that completes the family.}  Each kind of observation
constrains who may hold it.  \textbf{WF} does it for \frsh{}, \textbf{WUNLK}
for \unlk{} and \frbl{}, and this for \itr{}; the published set had exactly one
of the three.

It is \textsc{SyncStart}'s.  \textsc{SyncStart} puts every unlinked node on the
free list with the current readers as its bounding set, and \textbf{IFL} then
demands that every thread observing such a node as an \itr{} is in that set ---
it is the obligation that a grace period waits for everyone who could still be
looking.  \textbf{WULK} rules out the writer.  Nothing ruled out a thread that
is neither the writer nor a current reader, and against such a thread the grace
period would complete while a live reference remained, which is precisely the
accident the structure exists to prevent.

So of the three this is the one whose absence is not merely a failure of
locality.  \textbf{FNR}'s missing conjunct and \textbf{WUNLK} hold of every
reachable state and fail only to be provable locally; this one is what makes
reclamation safe, and it was unstated.  It is nonetheless as cheap as the other
two: every action preserves it, an \itr{} being granted only to the writer by
the two linking rules and only to a reader by a read, and \textsc{ReadEnd}
dropping a departing thread's observations in the same step as its readership.

\paragraph{19.\ UNQR --- Unique Reachability.\hfill\textsc{critical}}
\[
\textbf{UNQR} \defn \forall \rho,\rho',o \ldotp
  h^{*}(\sigma.rt,\rho) = \lfloor o \rfloor \land h^{*}(\sigma.rt,\rho') = \lfloor o \rfloor
  \implies \rho = \rho'
\]
\emph{Changed:} the original states
$h^{*}(\sigma.rt,\rho) \neq h^{*}(\sigma.rt,\rho') \implies \rho \neq \rho'$,
which is a theorem of equality --- if $\rho = \rho'$ then the two applications
of $h^{*}$ are equal by congruence, so the implication holds for \emph{any}
function whatsoever, in \emph{any} state, including states whose heap is a
cycle.  As published, \textbf{UNQR} constrains nothing.  The intended statement
is the converse: distinct paths from the root reach distinct nodes, i.e.\ the
structure is a tree.

\emph{Why this is critical:} \textbf{UNQR} is the tree-shape invariant on which
every framing side condition in \textsc{T-UnlinkH}, \textsc{T-Replace} and
\textsc{T-Insert} depends --- it is what licenses concluding node disequality
$o_x \neq o_y \neq o_z$ from path disequality $\rho \neq \rho' \neq \rho''$.
Six lemmas contain a \textbf{UNQR} case; all six discharge it against, or appeal
to it as, a statement with no content.  These are not restatement fixes.  They
have to be redone.

\subsection{Impact on the existing proof cases}
\label{sec:inv-impact}

The semantic changes above invalidate specific cases in Section
\ref{lem:lematom}.  Ordered by how much work they imply.

\begin{enumerate}
\item \textbf{UNQR (\S17)} --- \textsc{Unlink}, \textsc{Replace},
  \textsc{Insert}, \textsc{ReadHeap}, \textsc{WriteFreshField},
  \textsc{Alloc}.  Every \textbf{UNQR} case, plus every case that cites
  \texttt{.UNQR} as a hypothesis --- notably the \textbf{ULKR} cases of
  \textsc{Unlink} and \textsc{Replace}, which argue ``if $o'$ were observed as
  \itr{} that would conflict with \textbf{UNQR}''.

  \emph{Done for the mutations.}  The three heap-mutating cases are discharged
  against \S\ref{sec:inv-iris}'s dichotomy rather than by the local
  path-distinctness arguments the report gives, which proved something strictly
  weaker than the invariant.  \textsc{ReadHeap} mutates no edge, and
  \textsc{WriteFreshField} and \textsc{Alloc} write only an unreachable node, so
  all three are the unconditional case.  What is left of this item is its tail:
  the \textbf{ULKR} cases that cite \texttt{.UNQR}, which turn on the
  observation map rather than on reachability.

\item \textbf{FR (\S8)} --- \emph{done.}  The corrected conjunctive form does
  make these easier, as expected, but neither is the ``trivial'' the report
  claims: the action \emph{does} create an edge to the fresh node, and the case
  turns on the observation changing from \frsh{} to \itr{} in the same step, so
  that \textbf{FR}'s hypothesis stops applying to it.

\item \textbf{FPI (\S11)} --- \emph{done, after repairing two rules.}  Expected
  to be among the easiest items, this was instead the sixth defect.  \textsc{T-Replace} and \textsc{T-UnlinkH} both make a node
  \unlk, which by \textbf{WULK} costs it its \itr{} observation; if a
  \textsf{rcuFresh} variable in $\Gamma$ has a field pointing at that node,
  \textbf{FPI} holds before the write and fails after.  The rules' aliasing
  premises quantify only over \textsf{rcuItr} variables, and $\llbracket \Gamma
  \rrbracket$ is an intersection rather than a separating conjunction, so
  nothing excludes it; \textsc{T-WriteFH} is what makes the state reachable.
  The failure is not confined to \textbf{FPI}: $\llbracket
  \textsf{rcuFresh}\,\mathcal{N}\rrbracket$ carries the same requirement, so the
  affected lemmas fail to establish the denotation of their own
  post-environment.  Both rules now carry the repair --- a premise excluding
  \textsf{rcuFresh} predecessors of the node being unlinked, $y \neq o$ for
  \textsc{T-Replace} and $m' \neq z$ for \textsc{T-UnlinkH} --- and with it both
  cases are immediate.  \textsc{Insert} needs nothing, since it unlinks nothing.
  The binary search tree deletion still type checks: the fresh node it builds
  points at the children of the node being replaced, never at that node.  The
  counterexample showing the premise is not redundant is machine checked as
  \texttt{FPI\_not\_preserved\_by\_replace}.

\item \textbf{RITR (\S13)} --- \emph{done.}  The composition step of
  \textsc{WriteFreshField} now uses the corrected statement in the only form
  that carries information: not ``every reader observes every location as
  \itr'', which lets one conclude anything, but ``a reader observes nothing as
  \unlk, \frbl{} or \frsh'', which with \textbf{WNR} puts the written node
  outside the readers' fragment.

\item \textbf{HD (\S15)} --- \emph{done, after correcting the invariant.}  This
  item was expected to be the easiest --- four empty cases and one dismissed as
  trivial --- and the dismissed one turned out to be a defect: the
  published \textbf{HD} is false, and \texttt{Free} falsifies it.  See \S15 for
  the counterexample and the repair.  With the source restricted to non-detached
  nodes the \textsc{Free} case goes through by exactly the argument this list
  predicted, with \textbf{ULKR} supplying it
  (\texttt{free\_preserves\_HD}); under the published statement it cannot be
  proved at all, which is presumably why it was skipped.

  The empty cases were \textsc{Replace},
  \textsc{ReadHeap}, \textsc{WriteFreshField} and also \textsc{Alloc}, which this
  list previously missed; \textsc{Free} was worse than empty, folding
  \textbf{HD} into a blanket claim that all seventeen axioms are trivial.  Each
  is short, as expected: \textsc{ReadHeap} writes no field at all, the two field
  writes store a location the precondition already types as
  $\textsf{rcuItr}$ or $\textsf{rcuFresh}$ and hence already places in $dom(h)$,
  and \textsc{Alloc} only enlarges the domain with a node whose fields are all
  \textsf{null}.

  \textsc{Free} is the one that is not trivial, and it was the one dismissed as
  such: it is the only action that \emph{shrinks} $dom(h)$, so it must establish
  that no surviving edge points at the freed location.  That obligation is real
  rather than an artefact --- a state with a dangling predecessor satisfies
  \textbf{HD} before the free and violates it after
  (\texttt{HD\_h\_free\_needs\_no\_incoming}) --- and it is discharged from the
  transitive \textbf{ULKR} together with \textbf{FLR} and \textbf{IFL}.  The
  four preservation arguments are machine checked as \texttt{HD\_h\_upd},
  \texttt{HD\_h\_alloc} and \texttt{HD\_h\_free} in \texttt{rocq/HeapPaths.v}.

\item \textbf{ULKR (\S5)} --- the \textbf{ULKR} cases themselves are unaffected,
  since they establish the one-step form and the closed hypothesis is only
  harder in the antecedent.  What must be revisited is Section
  \ref{sec:lemmas}'s claim that \textbf{ULKR} and \textbf{FLR} together keep
  unlinked nodes unreachable from the root: that is the transitive form, and it
  needs the corrected invariant.  Both halves of it are now machine checked ---
  \texttt{ULKR\_closed\_reaches} and \texttt{FLR\_chains} --- so what remains at
  that citation site is to state the consequence in terms of them rather than of
  the one-step formulas.

\item \textbf{WULK, IFL, FLR, RINFL, UNQRT} --- restatement only.  This was
  conditional on the \textbf{FLR} direction; that question is now settled in
  favour of the reading already used in the proofs
  (\S\ref{sec:inv-open}), so the condition is discharged.  Recheck each citation
  site rather than assuming.
\end{enumerate}

\paragraph{Machine-checked status.}  Nine of the defects above are mechanized
in \texttt{rocq/WellFormed.v}: \texttt{UNQR\_orig\_vacuous} and
\texttt{UNQR\_orig\_admits\_non\_trees};
\texttt{RITR\_orig\_contradicts\_fresh};
\texttt{FR\_orig\_admits\_writer\_local\_alias};
\texttt{FPI\_orig\_contradicts\_FNR}; \texttt{ULKR\_does\_not\_chain}
together with its repair \texttt{ULKR\_closed\_reaches};
\texttt{HD\_not\_preserved\_by\_free}; and
\texttt{FNR\_pub\_admits\_freeable\_fresh}, which exhibits a state satisfying
the other eighteen invariants and the published \textbf{FNR} in which a node
is both \frsh{} and \frbl{} --- with \texttt{link\_freeable\_breaks\_ULKR} in
\texttt{rocq/Actions.v} showing that linking such a node breaks \textbf{ULKR},
and \texttt{link\_fresh\_ULKR} showing that the repaired \textbf{FNR} closes
the case; and
\texttt{unlinked\_observations\_need\_not\_be\_the\_writers}, which does the same
for \textbf{WUNLK} --- a state satisfying the other eighteen in which a node is
unlinked by a thread that is not the one holding the lock, so that the
\textbf{ULKR} case of the two unlinking rules cannot close; and
\texttt{iterators\_need\_not\_be\_active}, the same for \textbf{WITR} --- a state
satisfying the other eighteen in which a node is observed as an iterator by a
thread that is neither the writer nor a reader, against which
\textsc{SyncStart} would open a grace period that does not wait for it.  A
further one,
\texttt{FPI\_not\_preserved\_by\_replace}, is a defect in
\textsc{T-Replace} and \textsc{T-UnlinkH} rather than in an invariant
statement; both rules now carry the repairing premise.  Nine more are described
where they were found, and it is worth grouping them rather than numbering them
in sequence, because they are not all of a kind.

Two are defects in the model rather than in an invariant: \textsf{WellFormed} is
not closed under the published composition, so the framing the soundness
argument needs does not hold; and allocation makes the heap infinite, so no
points-to assertion exists over it.  Both are in \S\ref{sec:inv-mech}.

The other seven are further invariant defects, and they arrived in three waves.
Giving the linking rules a resource reading found three: the free list's
unconstrained domain (\S\ref{sec:defect-fld}) and two observations that
constrain nobody (\S\ref{sec:defect-obs}).  Giving the \emph{read side} one
found two more, at its two ends: \S\ref{sec:defect-br}, the bounding set is
unconstrained, so nothing makes the grace period terminate; and
\S\ref{sec:defect-wunlkw}, \textbf{WUNLK} is guarded by the lock, so it says
nothing in the state a reader starts in.  Attempting the reader's own read found
the last two: \S\ref{sec:defect-samesnap}, free-list entries need not agree,
which \textsc{SyncStart} carries as a hypothesis; and \S\ref{sec:defect-frw},
fresh-reachability is guarded by the stack, which is \textbf{WFresh}'s defect
one invariant along.

That is twenty: sixteen invariants, two rule defects, two model defects.  The
second rule defect is the one \S\ref{sec:inv-mech} reaches by chaining the rules
into a program: the mutation rules assume every field is an RCU field, which no
system with scalar fields satisfies and which a class declaration forbids.  The remaining
\textsc{major} and \textsc{minor} items are restatements whose defects are
syntactic (capture, scope, unused binders) and are visible without proof.

\paragraph{An invariant defect: the free list's domain is unconstrained.}
\label{sec:defect-fld}
Three invariants mention the free list --- \textbf{IFL}, \textbf{FLR} and
\textbf{RINFL} --- and all three are conditional on an entry existing.  None of
them says where entries may exist.  Nothing in the nineteen therefore rules out
an entry at a node that is not detached, not reachable, and not even in the
heap.

This is not a curiosity, and it was found the way the other six were: by trying
to prove something.  \textsc{SyncStart}'s specification says the grace period
covers \emph{exactly} the detached nodes, and re-establishing ``exactly''
requires knowing that the entries already present were at detached nodes too.
It is the one hypothesis of the snapshot step that no invariant supplies.
\texttt{free\_list\_domain\_is\_unconstrained} is the witness: a state
satisfying all nineteen with a free-list entry at a location that carries no
observations at all.

The repair is a twentieth conjunct,
\[
  \textbf{FLD} \quad\equiv\quad
  F(o) \text{ defined} \;\Rightarrow\; \exists t \ldotp
  o \mapsto \unlk{}_t \;\vee\; o \mapsto \frbl{}_t,
\]
which \textsc{SyncStart} establishes and which every other action either
preserves trivially or --- \textsc{Free} and \textsc{ReadEnd} --- preserves
because it only shrinks the free list.  It is stated in the development and used
as an explicit hypothesis where it is needed, rather than folded into \textsf{WellFormed}:
adding a conjunct means reproving fifteen actions, and that is a change to make
deliberately rather than in passing.

\paragraph{Two invariant defects: observations that constrain nobody.}
\label{sec:defect-obs}
Found the same way as \textbf{FLD}, one rule further along.
\textsc{T-LinkF-Null} publishes a fresh node, and to do so it must know that no
thread other than the writer observes that node.  Two invariants ought to supply
that between them, and neither does.

The first is that \rt{} marks the root.  Nothing says so.  \rt{} is the
anonymous observation --- it carries no thread --- and the nineteen constrain
where every \emph{tagged} observation may sit while saying nothing about this
one, so a well-formed state may record the root observation at a location that
is not the root.  That also weakens \textbf{UNQRT-b}, whose conclusion is
``iterator or root'' and which is read as ``every reachable node is the
writer's, except the root itself''.  Under the published set the exception is
not tied to the root at all.

The second is that a \frsh{} observation belongs to the lock holder.
\textbf{WFresh} is stated with a stack reference in its hypothesis --- a thread
that \emph{names} a fresh node holds the lock --- so a fresh observation whose
variable has left scope constrains nobody.  That makes \textbf{WFresh} weaker
than its own caption, and weaker than \textbf{WUNLK} and \textbf{WITR}, the two
this work adds, which are the same statement for the other two kinds of
observation.  The family is right; the member the published set already had was
the weak one.

One witness does for both (\texttt{root\_observation\_is\_unpinned} and
\texttt{fresh\_observations\_need\_not\_be\_the\_writers}): a state satisfying
all nineteen in which the root observation sits at a non-root location and a
fresh observation belongs to a thread that does not hold the lock.

The second repair is the obvious conjunct, and it is carried.  The first began
as one --- \textbf{RTO}, ``the root observation is at \rt{}'' --- and it is the
only one of the twenty that was later \emph{removed} rather than kept, which is
worth a paragraph because the reason generalizes.

The root observation is anonymous by design: being the root is a property of
the structure, not of anyone looking.  The encoding nonetheless filed every
observation under an observing thread, so the root observation had to be filed
under \emph{some} thread, and then three things followed.  \textbf{RTO} had to
be added, because nothing tied the filed observation to \rt{}.  \textsc{ReadEnd}
could not simply drop the departing thread's entries, because the root
observation might be in one --- there is a state where it is, and dropping them
breaks \textbf{UNQRT-b}; so \textsc{ReadEnd} carried an extra hypothesis about
what it preserved.  And the invariant's clause relating a reader's entries to
its registration had to be stated about \emph{tagged} observations rather than
about entries, which made it too weak to decide whether a reader has an entry at
a given node --- which is exactly what the reader's read has to know.

Reading the root observation off \rt{} instead removes all three at once.
\textbf{RTO} becomes true by computation (\texttt{to\_LState\_t\_RTO}, whose
proof is the identity function); the state that used to be the counterexample is
no longer well formed as an encoding (\texttt{root\_has\_no\_owner}), and
\textsc{ReadEnd} deletes its keys outright; and the domain clause becomes a
statement about entries, which is what \texttt{read\_update} needs.  The witness
above still stands --- the published nineteen really do not imply \textbf{RTO},
so the defect was real --- but no action has to preserve it.  The pattern is
that an invariant can be an artefact of an encoding choice, and that the cheaper
repair at the time (add a conjunct) was the wrong half of the choice.

Sixteen of the twenty were found by attempting a proof rather than by reading ---
\textbf{ULKR}'s failure to chain, \textbf{FPI}'s failure to survive unlinking,
\textbf{HD}'s failure to survive \texttt{Free}, \textbf{FNR}'s failure to
support \textbf{ULKR} under linking, \textbf{WUNLK}'s absence, which is the same case
one step further along, and \textbf{WITR}'s, which is \textsc{SyncStart}'s, \textbf{FLD}'s, which is
\textsc{SyncStart}'s other one, the failure of closure under composition, which
is what attempting the framing property ran into, the infinite heap, which is
what choosing a physical-state predicate ran into, and the two unconstrained
observations, which are what the last linking rule ran into, and the
\textbf{ULKR}/\textbf{OW} disagreement, which is what chaining the rules into a
program ran into, and the unconstrained bounding set and \textbf{WUNLK}'s guard,
which are what giving the read side a resource reading ran into, and the
entries' having to agree and fresh-reachability's guard, which are what the
reader's read ran into --- and seven of the sixteen are cases the report calls
trivial.  That is now a pattern rather than a coincidence, and it is the
strongest thing the mechanization has to say about the original argument.

The corrected invariants are also satisfiable
(\texttt{WellFormed\_satisfiable}), so the repairs are not vacuously strong.
What the development establishes positively --- the reachability theory, the
denotations, and the atomic-action lemmas --- is collected in
\S\ref{sec:inv-mech}.

\paragraph{An invariant defect: the bounding set is unconstrained.}
\label{sec:defect-br}
Exactly one of the nineteen mentions $B$: \textbf{RINFL} says every thread in a
free-list entry is a bounding thread.  That bounds the entries above by $B$ and
says nothing whatever about $B$ itself.

$B$ is not decoration.  \textsc{SyncStop} blocks until it is empty, and
\textsc{ReadEnd} is the only action that takes a thread out of it --- and
\textsc{ReadEnd} is a \emph{reader's} action.  So a bounding thread that is not
a reader can never be removed, and the grace period never completes.  Nothing in
the nineteen rules that out.
\texttt{bounding\_threads\_need\_not\_be\_readers} is the witness: a state
satisfying all nineteen with a thread in $B$ that is not a reader and holds no
observation at all.

The repair is the conjunct \textsc{SyncStart} in fact establishes, since it sets
$B$ to the current readers:
\[
  \textbf{BR} \quad\equiv\quad \forall t.\; t \in B \;\Rightarrow\; t \in R .
\]
Every action preserves it.  \textsc{ReadEnd} narrows $R$ and $B$ together,
\textsc{SyncStop} empties $B$, and no other action touches either.  It is the
exact counterpart of \textbf{WITR}: that one says who may hold an observation,
this one says who may hold up a reclamation.

Two things make this worth stating rather than waving at.  The first is that
without it the writer may bound its own grace period: \textbf{WNR} says the
writer is not a reader, and with \textbf{BR} that gives
\texttt{writer\_not\_bounding} --- the thread running the wait is never one of
the threads it waits for.  Without \textbf{BR} nothing rules that out, and
\textsc{SyncStop} deadlocks against itself.

The second is that the obstruction is provable, not merely arguable.  This
development has no operational semantics and so cannot state liveness; what it
can state is that the obstruction survives every step that could clear it.
\texttt{grace\_period\_stuck} is that statement: for any sequence of
\textsc{ReadEnd}s, taken in any order by any threads that are readers, a
bounding thread that is not a reader is still in $B$ afterwards.  That is the
safety-shaped half of ``the grace period never completes'', and it is the half a
development without a semantics is entitled to.

\paragraph{An invariant defect: WUNLK is guarded out of usefulness.}
\label{sec:defect-wunlkw}
\textbf{WUNLK} says detaching observations are the writer's --- but only
\emph{while a writer holds the lock}.  In an unlocked state it says nothing, so
nothing rules out a stray \unlk{} or \frbl{} observation belonging to an
arbitrary thread.  \texttt{detaching\_observations\_survive\_the\_lock} is the
witness: a state satisfying all nineteen, with the lock free and a node observed
\unlk{} by a thread that is not and never was the writer.

This is the defect \textsc{ReadBegin} runs into, and the unlocked state is
precisely when a thread starts reading.  A thread entering a read-side critical
section must hold no detaching observation, or \textbf{RITR} fails the moment it
becomes a reader.  With the lock held that does follow, from \textbf{WUNLK} and
the entering thread not being the writer.  With the lock free nothing supplies
it, and the action lemma has carried it as a hypothesis for want of an invariant
that does.

The repair is \textbf{WFreshW}'s, transposed --- state it unguarded:
\[
  \textbf{WUNLKW} \quad\equiv\quad
  o \mapsto \unlk{}_t \;\vee\; o \mapsto \frbl{}_t
  \;\Rightarrow\; \mathit{lk} = t ,
\]
carried, like \textbf{FLD}, as an explicit conjunct of the Iris invariant rather
than folded into \textsf{WellFormed}.  It is true of every reachable state for the reason the development already
records at the other end: \textsc{WriteEnd}'s own premise is that the critical
section leaves nothing detached.  So the invariant is not new information ---
it is information the rules already depend on, written down where an invariant
can be used.

That \textbf{WFresh} and \textbf{WUNLK} fail in the same way, one invariant
apart, and that the revision had already repaired the first without noticing the
second, is the part worth reporting.  The guard is what makes each of them
provable locally; it is also what makes them useless to any rule that runs
outside a write critical section, which is every rule on the read side.

\paragraph{An invariant defect: free-list entries need not agree.}
\label{sec:defect-samesnap}
\textsc{SyncStart} takes one snapshot, so every entry it writes holds the same
set.  That is a hypothesis of the snapshot step --- it is what makes \textbf{FLR}
hold with equality rather than mere inclusion along a chain of detached nodes
--- and it is stated there as a condition on the action.  It is also what the
\emph{reader's} read needs, and there it has to be an invariant, which is how we
found it.

\textbf{IFL}'s obligation on a reader following an edge is that a reader taking
a reference to a node already on the free list is one of the threads its grace
period waits for.  The reader is at $o$ and follows an edge to $z$.  If $z$ is
on the free list, \textbf{FLR} says its predecessor $o$ is too, with a set
\emph{containing} $z$'s; \textbf{IFL} applied to $o$, which the reader observes
as an iterator, says the reader is in $o$'s set.  What is wanted is that it is in
$z$'s, and the inclusion runs the wrong way.  The two facts do not meet.

With the entries equal they meet at once, and \texttt{read\_bound} is four
lines.  Nothing in the nineteen forces them to agree:
\texttt{free\_list\_entries\_need\_not\_agree} is a state satisfying all of
them, and \textbf{BR} besides, with two entries holding different sets.  The
repair is the conjunct \textsc{SyncStart} already establishes --- in the epoch
reading it is the statement that every stamp is the same epoch, which is exactly
the fact \texttt{sync\_start\_update} proves on the way through.

This is the fourth of the same kind, and by now the pattern is worth naming: a
condition the published proof carries as a \emph{hypothesis of one action}, and
which some other rule needs as an \emph{invariant}.  \textbf{WITR} was
\textsc{SyncStart}'s, \textbf{FLD} was \textsc{SyncStart}'s other one,
\textbf{WUNLKW} was \textsc{WriteEnd}'s, and this is \textsc{SyncStart}'s third.
An invariant set assembled by asking what each action needs locally will miss
exactly these, because locally there is nothing to miss.

\paragraph{An invariant defect: fresh-reachability is guarded by the stack.}
\label{sec:defect-frw}
The fifth of the family, and the reader's read found it the way it found the
fourth --- by needing, as an invariant, something the published proof carries as
a hypothesis of one action.

\textsc{T-ReadH} gives a reader an \itr{} observation on the node it has just
reached.  If that node were \frsh{}, \textbf{FNR} would fail at once, since a
fresh node holds no iterator observation.  So the read needs to know the node it
reaches is not fresh, and what ought to supply that is \textbf{FR}: a fresh node
has no incoming edge, so a node reached by following one is not fresh.

\textbf{FR} does not supply it.  It is stated with a stack reference in its
hypothesis --- a thread that \emph{names} a fresh node --- exactly as
\textbf{WFresh} was before the revision unguarded it (\S\ref{sec:defect-obs}),
and it is unusable for the same reason: the observation outlives the variable.
A fresh node whose variable has been rebound constrains nobody, and nothing then
stops an edge pointing at it.  \texttt{fresh\_may\_be\_pointed\_at} is the
witness --- a state satisfying all nineteen in which a detached node points at a
fresh one.  The repair is \textbf{WFresh}'s, and it is true of every reachable
state because a node stops being fresh at the moment it is published.

That the revision unguarded \textbf{WFresh}, and then had to unguard
\textbf{WUNLK} (\S\ref{sec:defect-wunlkw}), and now \textbf{FR}, is the pattern
stated three times.  Each was provable locally \emph{because} of its guard, and
useless outside the critical section \emph{because} of it.

\subsection{Open questions}
\label{sec:inv-open}

\paragraph{Direction of the inclusion in FLR (\S6) --- resolved.}
The original concludes $\exists T' \subseteq T \ldotp F(o' \mapsto T')$: the
predecessor's reader set is a \emph{subset} of the successor's.  This was
initially left open, on the stated grounds that the temporal story argued
\emph{against} what the report writes.  That was an error --- the two agree, and
the direction is forced rather than a matter of adjudication.

A free-list entry is populated at \texttt{SyncStart} with the then-current
reader set and shrinks as readers execute \texttt{ReadEnd}.  Since
$\sigma.h(o',f') = o$ puts $o'$ above $o$, and \textbf{ULKR} propagates upward
--- an unlinked child forces an unlinked parent --- $o'$ must already be
unlinked at the instant $o$ becomes unlinked, so $o'$ is unlinked no later than
$o$.  Grace periods cannot overlap: \textsc{T-Sync} types
\texttt{SyncStart};\texttt{SyncStop} as a single compound statement and there is
one writer.  So there are only two cases.  If $o'$ and $o$ were unlinked in
different write critical sections, the earlier grace period has already
completed, $o'$ is \frbl{} and $F(o') = \emptyset$.  If in the same one, both
entries were populated at the same \texttt{SyncStart} from a single snapshot and
are equal.  Either way $T_{r}' \subseteq T_{r}$.

Non-overlapping grace periods are also what answers the obvious objection, that
entries created at different \texttt{SyncStart}s snapshot unrelated sets of
threads and so need not nest at all: two \emph{live} snapshots are never
compared, because the earlier one has drained to $\emptyset$ before the later
one is taken.

With the direction fixed, \textbf{FLR} chains, and the one-step form needs no
repair --- unlike \textbf{ULKR}, its hypothesis and conclusion are the same
predicate, so only the induction is required.  The result is
\texttt{FLR\_chains} in \texttt{rocq/WellFormed.v}: anything reaching a
free-list entry is itself one, with a contained set.  Together with
\texttt{ULKR\_closed\_reaches} this is what keeps a node awaiting reclamation
unreachable from the root, the root being on neither the free list nor the
unlinked set.  \texttt{WellFormed} is correspondingly no longer parameterised by
the inclusion.

\paragraph{Whether AWRT should be derived rather than assumed.}
See \S\ref{sec:inv-iris}.

\paragraph{$\mathcal{N}_1 = \mathcal{N}_2$ in \textsc{T-Replace} is weaker than the proof needs (resolved).}
Found while mechanizing the \textsc{T-Replace} case
(\texttt{rocq/HeapPaths.v}).  Preserving \textbf{UNQR} across the replacement
requires that the fresh node's subtree \emph{be} the old node's subtree, i.e.\
$\forall g \ldotp \sigma.h(n,g) = \sigma.h(o,g)$ on \RCU{} fields.  The rule
supplies $\mathcal{N}_1 = \mathcal{N}_2$, equality of the two \emph{field maps}.  That is not the
same thing: a field map records only fields that have been read or set, and
neither $\llbracket \textsf{rcuItr}\,\rho\,\mathcal{N} \rrbracket$ nor $\llbracket
\textsf{rcuFresh}\,\mathcal{N} \rrbracket$ constrains a field outside $dom(\mathcal{N})$ --- so
the two subtrees may differ on an untracked field while $\mathcal{N}_1 = \mathcal{N}_2$ holds.

\textsc{T-Insert} has the same \emph{symptom}, in the same place: its
$\forall_{f_2\in dom(\mathcal{N}_1)} \ldotp f_4 \neq f_2 \implies \mathcal{N}_1(f_2) =
\textsf{null}$ constrains only the tracked fields, so an untracked field of the
fresh node could point anywhere and give the inserted node a second child.

\paragraph{Resolution.}  One tempting reading should be discarded first.  Section
\ref{sec:tslbl}'s ``we presume all objects contain all fields'' is about object
\emph{layout} --- every object has every field present, so the system need not
do C-style field-presence checking --- and says nothing about $dom(\mathcal{N})$.  The
rules are explicit that it does not: \textsc{T-Alloc} produces
$\mathcal{N}_{\emptyset}$, \textsc{T-ReadH} produces $\mathcal{N}_{\emptyset}$ for the
variable it reads into, and \textsc{T-WriteFH} requires $f\notin dom(\mathcal{N}')$ and
returns $\mathcal{N}'([f\rightharpoonup z])$.  A field map is by construction ``the
fields read or written so far''; totality is contrary to the design, not an
unstated intent the system fails to enforce.

The two rules therefore need fixes in different places, because \textsc{T-Insert}'s
obligation is about the \emph{fresh} node and \textsc{T-Replace}'s is about the
\emph{pre-existing} one.

\textsc{T-Insert} is fixed in the denotation, with no change to the rule.
$\llbracket \textsf{rcuFresh}\,\mathcal{N} \rrbracket_{tid}$ now also requires every
\RCU{} field outside $dom(\mathcal{N})$ to be \textsf{null}.  This is an invariant of the
fresh-node lifecycle rather than a new proof obligation: \textsc{T-Alloc}
establishes it, and \textsc{T-WriteFH} --- the only rule that writes a field of a
\textsf{rcuFresh} object --- preserves it by moving exactly the written field into
$dom(\mathcal{N})$.  With it, the rule's existing null condition on $dom(\mathcal{N}_1)$ gives
\texttt{PointsOnlyAt} outright.

\textsc{T-Replace} cannot be fixed the same way.  \texttt{Mirrors} quantifies over
the fields of \emph{both} nodes, and the second is an \textsf{rcuItr} inside the
structure whose untracked fields are legitimately non-null, so ``unlisted implies
\textsf{null}'' would simply be false of it.  That one needs an explicit side
condition on the rule, $\forall_{f'} \ldotp \textsf{FType}(f')=\RCU \implies f'
\in dom(\mathcal{N}_1)$; together with $\mathcal{N}_1 = \mathcal{N}_2$ the two maps then pin the entire
\RCU{} footprint of both nodes and \texttt{Mirrors} follows from the denotations.

Only the \textsc{T-Replace} condition costs the checker anything, and in practice
it is bookkeeping: writing a field of the fresh node with \textsc{T-WriteFH}
already requires that field to be in the replaced node's map, so any replacement
that copies every \RCU{} field satisfies it by construction.  The binary search
tree deletion sets both \texttt{Left} and \texttt{Right} of the fresh node before
replacing, and so still type checks.

The mechanization continues to take \texttt{Mirrors} and \texttt{PointsOnlyAt} as
hypotheses --- \texttt{HeapPaths.v} models the heap, not the type system --- but
they are now supplied by the revised rule and denotation rather than assumed
away.

\subsection{The encoding: raw heap plus $h^{*}$}
\label{sec:inv-iris}

Four of these invariants --- \textbf{UNQR} (\S17), \textbf{UNQRT-a} (\S16),
\textbf{OW} (\S1) and \textbf{HD} (\S15) --- exist only because the model
represents the structure as a raw heap together with a reachability function
$h^{*}$, and must then \emph{assert} that the result is a tree.  Carrying an
inductive representation predicate instead would make all four derivable.

\paragraph{Decision.}  We keep the raw heap.  Two reasons.  First, fidelity: the
existing soundness argument is phrased over $\sigma.h$ and $h^{*}$, so its
structure --- once repaired --- transfers directly, whereas a representation
predicate would require re-deriving each case.  Second, the split a
representation predicate appears to buy is illusory here: unlinked-but-not-yet
-freed nodes sit \emph{outside} the tree while readers may still hold them, so
the invariant would have to be
$\textsf{tree\_rep}\;rt\;t \ast \big(\ast_{o \in dom(F)} \textsf{detached}\;o\big)$
and the detached region would need its own reachability reasoning regardless.

\paragraph{Consequence.}  Those four invariants must be re-established by hand
in every atomic-action case, and they are the fiddly ones --- \textbf{UNQR}
especially, since preserving ``distinct paths reach distinct nodes'' across a
field write is precisely the step the technical report's proofs gesture at.  The
mitigation is to prove the reachability theory once and discharge every action
against it, rather than re-deriving a path argument per rule.

That theory is \texttt{rocq/HeapPaths.v}.  Its central result is a dichotomy:
after a single field write, every path either behaves exactly as it did before,
or it factors through the written edge.
\[
\begin{array}{l}
\forall h, o_u, f_u, v, \rho, o \ldotp
  h^{*}[o_u.f_u \mapsto v](o,\rho) = h^{*}(o,\rho) \\
\qquad {} \lor\; \exists \rho_1,\rho_2 \ldotp \rho = \rho_1 \cdot f_u \cdot \rho_2
  \land h^{*}(o,\rho_1) = o_u
\end{array}
\]
Since each of \textsc{T-UnlinkH}, \textsc{T-Replace} and \textsc{T-Insert}
performs exactly one write, all three reduce to this plus a statement of how the
factored suffix is rewritten.  All three are proved:
\[
\begin{array}{lll}
\textsc{T-Replace}  & \text{length-preserving} & \texttt{UNQR\_replace} \\
\textsc{T-Insert}   & \text{lengthened by one} & \texttt{UNQR\_insert} \\
\textsc{T-UnlinkH}  & \text{shortened by one}  & \texttt{UNQR\_unlink}
\end{array}
\]
Two things worth recording from doing them.  \textsc{T-UnlinkH}, which the
technical report's proofs gesture at hardest, turns out to need \emph{fewer}
hypotheses than the other two: with no fresh node there is nothing to mirror and
no fields-are-null condition, and the rule's
$\forall_{f\in dom(\mathcal{N}_1)} \ldotp f \neq f_2 \implies \mathcal{N}_1(f) = \textsf{null}$
plays no part --- it bounds how much of the structure becomes unreachable, which
is \textbf{ULKR}'s business, not \textbf{UNQR}'s.  And the case that has to be
excluded in \textsc{T-Insert} and \textsc{T-UnlinkH}, but not in
\textsc{T-Replace}, is a collision between an \emph{unchanged} path and a
rewritten one; ruling it out is why the characterisation must record that the
unchanged path does not traverse the written edge.

Both \textsc{T-Insert} and \textsc{T-UnlinkH} need a side condition saying that
nothing below the write can reach back above it.  Neither is an extra proof
obligation: both follow from \textbf{UNQR} itself together with the reachability
of the written node, which is what the rules supply, and the mechanization
discharges them (\texttt{UNQR\_insert\_reachable},
\texttt{UNQR\_unlink\_reachable}) so that the composition is checked rather than
asserted.

The residue --- the observation map $O$, the free list $F$, the grace-period
invariants (\textbf{IFL}, \textbf{RINFL}, \textbf{FLR}) and the fresh-node
invariants (\textbf{FR}, \textbf{WF}, \textbf{FNR}, \textbf{FPI}) --- is
genuinely about ghost state and becomes ghost state in the Iris encoding.

\subsection{The ghost state, and what it had to be}
\label{sec:inv-ghost}

$O$ becomes a finite map from locations to sets of observations and $F$ a map to
sets of thread identifiers, each under an authoritative resource algebra, so that
the invariant
$\textsf{inv}\;\mathcal{N}\;(\exists s \ldotp \textsf{phys} \ast \textsf{ghost}
\ast \ulcorner \textsf{WellFormed}\;s \urcorner)$
holds the authority while a thread holds an entry.  The pure layer is reused
unchanged: a function reconstructs a logical state from the ghost maps, and every
result above is stated over that reconstruction, so nothing has to be restated.

One choice there is worth recording, because getting it wrong is easy and the
symptom is remote from the cause.  The natural first reading of an observation
map is a resource algebra whose operation is \emph{union}: a thread owns some
observations, the authority owns at least those, and ownership is freely
duplicable.  That is wrong, and the reason is \textbf{WULK}.  Unlinking must
\emph{withdraw} an observation --- the replaced node loses \itr{} and gains
\unlk, and \textbf{WULK} makes the two exclusive --- whereas under a union
algebra every fragment is duplicable and therefore permanent, so an observation
once handed out could never be given up and no unlinking step could re-establish
the invariant.  Exclusive control of a location's entry is what the type system
actually wants for the writer, and it is what we use.  The union reading is
correct for a quantity that only grows; observations are not one.

\subsection{Status of the mechanization}
\label{sec:inv-mech}

Beyond the defects, the development discharges the following, all axiom-free.

The reachability theory of \S\ref{sec:inv-iris}, and heap-domain closure for
every action that touches the heap --- including \texttt{Free}, the only action
that removes a location and so the only one where the obligation is real; a
counterexample shows its side condition cannot be dropped.

The type denotations of Figure \ref{fig:denotingtypeenviroment}, and with them
the two hypotheses the reachability theorems take.  That the fresh node mirrors
the replaced one, and that its only outgoing \RCU{} edge is the one the rule
names, are \emph{derived} from the denotations together with the rules'
premises rather than assumed.  Both derivations depend on a repair: the first on
the side condition added to \textsc{T-Replace}, the second on the clause added
to $\llbracket \textsf{rcuFresh}\,\mathcal{N} \rrbracket$.  Neither proof closes
without it, so the repairs are load-bearing rather than bookkeeping.

Atomic-action lemmas for every action.  For the three heap mutations: each
preserves \textbf{UNQR} and \textbf{UNQRT}$_b$, and a field write preserves
\textbf{OW}, \textbf{HD}, \textbf{ULKR}, \textbf{FLR}, \textbf{FR} and
\textbf{UNQRT}$_a$ under a side condition apiece, each of which the rules
already supply --- with the one exception that produced the \textbf{FNR} repair
above.  Unlinking and replacement preserve \textbf{FPI}.  The \textbf{FPI} case
is the one whose failure forced the repair to \textsc{T-Replace} and
\textsc{T-UnlinkH}, and the
repaired premise is used exactly once in that proof, at the step that would otherwise
fail --- so that defect is closed from both ends, by a counterexample against the
published rule and a proof for the repaired one.  The argument does not depend on
which rule caused the write, so it is proved once and both rules are instances,
which is itself why the two needed the same repair.

For the reclamation protocol, the chain closes end to end.  Acquiring a
reference obliges the acquiring thread to bound any pending reclamation of the
node, which is \textbf{IFL}; \textsc{SyncStart} takes one snapshot of the
readers, so all entries it creates agree and \textbf{FLR} holds with equality
--- the same-critical-section case of the argument that fixed its direction;
\textsc{SyncStop} leaves the bounding set empty, so by \textbf{RINFL} every
entry is empty, which is the conjunct the \frbl{} denotation asks for;
\textsc{Free} then needs that nothing \emph{live} points at the node, and
\textbf{ULKR} supplies exactly that.  \textsc{ReadEnd} closes the loop: it must
drop a thread's observations, its bounding-thread membership and its free-list
entries \emph{together}, and the \textbf{IFL} and \textbf{RINFL} cases are
what show that none of the three can be taken alone.

\textsc{Free} is also the first action taken \emph{whole}: not one invariant at
a time but \textbf{WellFormed} itself, as \texttt{free\_preserves\_WellFormed}.
It is the right one to do first, because it needs no type denotations --- it is
a state transition rather than a typing rule --- and because it is the action
the report handled worst.  The result is lopsided in a way worth recording.
Eighteen of the nineteen fall out of two one-line facts, that freeing removes
edges and therefore removes paths: every invariant with an \textbf{Edge} or
\textbf{Reaches} hypothesis gets \emph{weaker}, and the eleven that do not
mention the heap are untouched outright.  \textbf{HD} is the only one that asserts
something must be \emph{in} the heap, and the only one with any content --- and it needs the
corrected statement, since under the published \textbf{HD} the theorem is false.
So the report's verdict on the case was right about eighteen nineteenths of it
and wrong about the part that mattered, which is a fair summary of what
"trivial" was doing throughout.

\textsc{SyncStop} follows, as \texttt{sync\_stop\_preserves\_WellFormed}.  It
does two things --- it returns once the bounding set is empty, and it turns
every \unlk{} observation into the matching \frbl{} one --- and the two halves
carry disjoint sets of invariants.  The observation change is stated by what it
does to observations rather than as a map operation, which is the choice
\textsc{ReadEnd} already makes and for the same reason: two properties suffice,
and they are the two directions.  Nothing appears that was not there, except
\frbl{} arriving from \unlk{}; and everything survives, \unlk{} as \frbl{}.
Every one of the nineteen cases is one of those two applied once.

Two have content.  RINFL is where the empty bounding set is spent: it says
every thread in a free-list entry bounds a reclamation, and afterwards there
are no bounding threads, so no entry has a thread in it --- which is exactly the
conjunct the \frbl{} denotation asks for, obtained rather than assumed.  And
FNR is \texttt{sync\_stop\_preserves\_FNR}: the conjunct added to FNR to make
\textsc{T-Insert} and \textsc{T-Replace} preserve \textbf{ULKR} is discharged
here from the conjunct FNR already had.  The repair therefore costs nothing at
the one action that could have been asked to pay for it, which is the sense in
which it was the right repair and not merely a sufficient one.

Because that step is stated abstractly, it is also checked to be
implementable: \texttt{sync\_stop\_WellFormed} instantiates both properties
with the obvious map on the observation state.  The check is not ceremony ---
\textbf{RITR} is the standing reminder that an abstractly stated condition can
be one no concrete step satisfies, and every theorem over it then holds for
nothing.

\textsc{ReadEnd} is the third, as \texttt{read\_end\_preserves\_WellFormed},
and it is the one where doing the action whole \emph{changed} the statement
rather than confirming it.  Two things had to be added that no per-invariant
case needed, and neither is bookkeeping.

The first is the departing thread's variables.  \textbf{RWOW} says a live
variable's referent carries an observation, and \textsc{ReadEnd} drops the
thread's observations; unless the thread's variables stop being live in the
same step, \textbf{RWOW} fails at once.  The type system does supply this ---
\textsc{ToRCURead} types the read critical section as a block, so the reader's
\textsf{rcuItr} variables are scoped to it and never reach the outer
environment --- but that is a fact about the rule, and it has to be carried into
the action's statement to be usable.  The earlier \textbf{IFL} and
\textbf{RINFL} cases did not need it because neither invariant mentions the
stack, which is exactly the kind of thing per-invariant work cannot notice.

The second is \rt{}, and it is an encoding mismatch rather than a defect in
anything published.  The observation map is keyed by observer, so that a thread
can own its own fragment; \rt{} has no observer --- change C0 keeps it
anonymous, being the root being a property of the structure rather than of
anyone looking --- so \textbf{ObsWF} permits it in any thread's entry, a
reader's included.  Describe \textsc{ReadEnd} as "the departing thread's entries
are gone", which is what the per-invariant lemmas use and what is correct for
every invariant about a thread's observations, and that description destroys
the root observation; \textbf{UNQRT}$_b$, the one invariant that asks about
\rt{}, then fails.  \texttt{read\_end\_naive\_breaks\_UNQRT\_b} is a well-formed
state where it does.  Either the encoding gives \rt{} a home no
\textsc{ReadEnd} can remove, or the action is stated so as to keep it; the
development takes the second, which costs one hypothesis against a change to
\textbf{ObsWF} that every lemma in the ghost-state layer depends on.

With those, the step is three properties of the observations, the same shape
\textsc{SyncStop} takes.  The pivotal hypothesis is that \textsc{ReadEnd} is a
\emph{reader's}: \textbf{RITR} then says the departing thread held no
detaching observation, so nothing any detachment invariant depends on is among
what is dropped.  Without \textbf{RITR} the action is unsound, and this is the
only place in the development where \textbf{RITR} is needed for anything but
its own preservation --- which is worth recording, given that \textbf{RITR} as
published is unsatisfiable in exactly the states where it is needed.
\textbf{WNR} is needed here too: the writer is not the departing thread, so its
\itr{} observations are among those that survive, which is what \textbf{FPI}
and \textbf{UNQRT}$_b$ need.  Unlike \textbf{RITR} it is needed in several
places besides, being what says the writer is not a reader --- the \textbf{RITR}
case of \textsc{T-UnlinkH}, \textsc{T-Replace} and \textsc{T-Alloc}, and the
\textbf{WULK} case of a reader's acquisition, all turn on it.

Finally the three heap mutations, and then every remaining action:
\textsc{ReadBegin}, the acquisition of a reference inside a read section,
\textsc{SyncStart}, \textsc{T-Alloc}, \textsc{T-WriteFH}, the three rules that
bind a writer's variable --- \textsc{T-Root}, \textsc{T-ReadS} and
\textsc{T-ReadH}, which are one lemma, differing only in where the node comes
from and not at all in what they do to the state --- \textsc{T-LinkF-Null}, and
\textsc{WriteBegin} and \textsc{WriteEnd}.  Fifteen theorems for the whole
system: the six RCU primitives and every type rule that changes the state, with
the three binding rules sharing one and a reader's read composing two.  The
structural and subtyping rules change no state and carry no action lemma.

\textsc{T-LinkF-Null} is the variant of \textsc{T-Insert} in which the fresh
node's \RCU{} fields are all null --- appending at the end of a list, or
inserting a leaf.  The paper mentions it and does not show it, and it is not an
instance of \textsc{T-Insert}: that rule's reachability argument requires the
fresh node to have exactly one outgoing edge, and here it has none.  It is the
easier of the two in exactly the place where \textsc{T-Insert} was hard, the OW
and HD cases that had to be re-discharged there being vacuous when the node is
the source of no edge.

\textsc{SyncStart} is where \textbf{WITR} is spent and the reason it exists, and
\textsc{T-Alloc} is Free's opposite: Free is the only action that removes a
location and so the only one that can strand an edge, while allocation extends
the domain and can strand nothing.  It is also where the fresh-node conditions
\textsc{T-Insert} and \textsc{T-Replace} consume --- nothing points at the node,
no other reference names it, its fields are null --- are established rather than
assumed.

\textsc{WriteBegin} is the one that does not go through as written, and it is
worth stating separately because the obstacle is an invariant rather than a
proof.  \textsc{RCU-WBegin} takes the lock and changes nothing else.  But
\textbf{UNQRT-b} says that while a writer holds the lock every reachable node
carries \emph{that writer's} \itr{} observation, and before the step there is no
writer, so the invariant said nothing.  A thread that has just taken the lock
holds no references, so on any structure larger than the bare root
\textbf{UNQRT-b} is false immediately after a lock acquire that only takes the
lock.

The step therefore has to grant the incoming writer an \itr{} observation on
every reachable node.  That is harmless in itself --- the observation map is
ghost state, and a writer inside its critical section may indeed reach anything
from the root, so recording it asserts nothing new about the heap --- but it is
not nothing, and neither the semantics nor the invariant discussion says it
happens.  Whether the right repair is that grant or a weaker \textbf{UNQRT-b} is
a question about the intended reading of the observation map, and we do not
settle it here.  What the mechanization establishes is that one of the two is
needed: the invariant as published and the action as published cannot both
stand.  It is the only place in the development where we record an obligation
without choosing how to discharge it.

\textsc{WriteEnd} is the counterpart, and unremarkable beside it: the departing
writer drops its observations and its variables in the same step as the lock,
exactly as \textsc{ReadEnd} does, and \textsc{ToRCUWrite}'s condition that
nothing is left \unlk{} or \frbl{} --- and, by the same argument, nothing left
\frsh{} --- makes most of the nineteen vacuous rather than merely easier.  That
is the sense in which a write critical section is a closed unit.  They are harder than the other three,
for a reason that appears only when the action is attempted in full: each
writes a field \emph{and} changes the writer's observations, and for
\textsc{Free}, \textsc{SyncStop} and \textsc{ReadEnd} the eleven heap-free
invariants were identities precisely because nothing touched the observation
map.  Here they are not.  \textsc{T-Insert} promotes a node from \frsh{} to
\itr{}, \textsc{T-UnlinkH} demotes one from \itr{} to \unlk{}, and
\textsc{T-Replace} does both at once, so ten of the eleven --- \textbf{RWOW},
\textbf{AWRT}, \textbf{IFL}, \textbf{WULK}, \textbf{WF}, \textbf{FNR},
\textbf{RITR}, \textbf{RINFL}, \textbf{WUNLK} and \textbf{WITR} --- all acquire
content.  \textbf{WNR} is the one that does not, being the only one of the
eleven that mentions no observation.

The observation change is again given by properties rather than by a map
operation, and here that is not merely for uniformity.  The revocation an
unlink performs is of \emph{the writer's} iterator observation and of nothing
else --- the readers still traversing the node keep theirs, which is the entire
point of the structure --- so an update replacing a location's whole
observation set would be wrong, and wrong in a way \textbf{FPI} cannot see and
\textbf{RWOW} can.  That is worth recording, because the per-invariant
\textbf{FPI} lemma uses exactly such an update.  It is sound there, its
conclusion never asking about a reader's observation, and it would not have
been here.

Two earlier repairs are spent in these three.  The premise excluding a
\textsf{rcuFresh} predecessor of the node being unlinked --- the repair
\textbf{FPI} forced --- is used twice in each unlinking rule, once in
\textbf{ULKR} and once in \textbf{FPI}.  \textbf{FNR}'s third conjunct is used
again in the \textbf{WULK} case of each of the three rules that link a fresh
node --- \textsc{T-Insert}, \textsc{T-Replace} and \textsc{T-LinkF-Null} ---
where the node just promoted to \itr{} has to be shown to carry neither \unlk{}
nor \frbl{}.  And \textbf{WUNLK} is the one thing these three needed that the
published set did not have.

One structural observation came out of doing the breadth.  Eleven of the nineteen
invariants --- \textbf{RWOW}, \textbf{AWRT}, \textbf{IFL}, \textbf{WULK},
\textbf{WF}, \textbf{FNR}, \textbf{RITR}, \textbf{RINFL}, \textbf{WNR},
\textbf{WUNLK}, \textbf{WITR} --- do not mention the heap, so a field write
transfers them definitionally and their cases are the identity.  The eight that do --- \textbf{OW}, \textbf{ULKR},
\textbf{FLR}, \textbf{FPI}, \textbf{FR}, \textbf{HD} and the two halves of
\textbf{UNQRT} --- each needed a real argument, and three of them turned up
defects.  The split falls exactly on that line.  It is worth stating because the
difference between the two kinds of trivial case is precisely what the original
proofs did not record, and that is where the defects hid.

Within those eight there is a second, finer split, and it is what located the
\textbf{FNR} defect.  Six --- \textbf{OW}, \textbf{ULKR}, \textbf{FLR},
\textbf{FR}, \textbf{HD}, \textbf{UNQRT}$_a$ --- are edge-local: a write only
ever \emph{removes} the edge it overwrites, and none of the six is endangered
by an edge going away, so each reduces to a single obligation about the edge
created.  \textbf{FPI} is not, since it needs the observation the same step
withdraws; and \textbf{UNQRT}$_b$ is not a write lemma at all, its hypothesis
being reachability, so it is discharged against the path characterizations
instead, one rule at a time.

Reducing the six that way is what turned a proof case into a proposition, and
it is the reason the \textbf{FNR} gap became visible: \textbf{ULKR}'s
obligation is a formula about $o_n$ that has to be discharged from the
invariants, and once written down it plainly could not be.  The defect it
exposed is not in any of the eight --- \textbf{FNR} is one of the eleven that a
write transfers definitionally --- which is the point.  An invariant can be
heap-free and still be the thing a heap obligation needs.

\textbf{UNQRT}$_b$ is worth a paragraph of its own, because it is the only
invariant that couples \emph{reachability} with the observation map --- \textbf{OW}
and \textbf{HD} couple observations with individual edges, which a write changes
locally, whereas \textbf{UNQRT}$_b$ asks about every path from the root --- and because
doing it exposes an asymmetry the report's case structure does not have.
\textsc{T-Insert} is easy: nothing loses an observation, and everything
reachable afterwards was either reachable before or is the inserted node, which
the step observes as an iterator.  \textsc{T-UnlinkH} and \textsc{T-Replace}
each \emph{revoke} an iterator observation, so for them the invariant needs
something \textbf{UNQR} does not say: that the node losing it has left the
structure.  \textbf{UNQR} says the result is still a tree, not which nodes are
in it, and this is the one case where the difference bites.  The two missing
facts are proved as \texttt{unlink\_unreachable} and
\texttt{replace\_unreachable} in \texttt{rocq/HeapPaths.v}, each by the same
path-length argument as \texttt{UNQR\_no\_back}; their proofs differ only
because the replacement characterisation, unlike the unlinking one, rewrites
paths through the written edge rather than excluding them.

What remains is worth stating precisely, because the shape of these theorems
invites overreading them.  Each says that an action preserves
\textbf{WellFormed}: the global invariants survive every atomic step.  Axiom
soundness asks for more.  In the form the soundness chapter states it,
\[\forall m \ldotp \llbracket \alpha \rrbracket (\lfloor \llbracket \Gamma_{1}
\rrbracket_{tid} * \{m\} \rfloor) \subseteq \lfloor \llbracket \Gamma_{2}
\rrbracket_{tid} * \mathcal{R}(\{m\}) \rfloor\]
the post-state must satisfy the denotation of the \emph{post-type environment},
not merely the invariants --- that after \textsc{T-Replace} the fresh reference
really is an \textsf{rcuItr} at the replaced node's path and the replaced one
really is \textsf{unlinked}, with their field maps as the rule records them.

That is now proved, for every component of every writer-side rule's output
environment: fifteen theorems.  \texttt{alloc\_post\_env} gives
$x : \textsf{rcuFresh}\,\N_\emptyset$ and \texttt{free\_post\_env} gives
$x : \textsf{undef}$; \texttt{bind\_post\_env} gives $y :
\textsf{rcuItr}\,\rho\,\N$ for \textsc{T-Root}, \textsc{T-ReadS} and
\textsc{T-ReadH}; \texttt{write\_fresh\_post\_env} gives the fresh node's
extended field map.  Each of the four heap mutations has all of its components:
the node retyped, the \emph{parent} with the written field recorded in its map,
and --- where the write moves a node --- the third reference whose path changes.
\textsc{T-Insert} lengthens the displaced node's path by one field
(\texttt{insert\_post\_env\_displaced}) and \textsc{T-UnlinkH} shortens the
promoted node's by one (\texttt{unlink\_post\_env\_promoted}), which are the two
places where the type system's paths and the heap's have to be shown to agree
after the structure moves.

Two of those cases need an invariant rather than a premise.  The parent's map,
in \textsc{T-UnlinkH} and \textsc{T-Replace}, may not already record a field
pointing at the node the rule unlinks, and it is OW that rules this out --- the
parent and the unlinked node are both the writer's iterators, hence neither
detached, so a second edge between them is impossible.

What is \emph{not} proved is the read-side.  A reader's \textsf{rcuItr} carries
no path and no field map, and that denotation is not in this development at
all, so the reader's instances of \textsc{T-ReadS} and \textsc{T-ReadH} have
their invariant half and not their type half.  Writing it down is a small piece
of work that has to precede the proof rather than being part of it.

\paragraph{Framing, and a defect in the framework's requirement.}  The Views framework requires the
views to form a partial commutative monoid and the validity predicate to be
closed under the operation; framing is unsound without it.  Section
\ref{sec:soundness} gives the operation --- Figure \ref{fig:comp} --- and asserts
the closure: \emph{WellFormed states are closed under both composition (with
another WellFormed state) and interference}.

That assertion does not hold, and the reason is structural rather than
incidental: the heap component of the composition is not separating.
$\bullet_h$ is defined at \emph{every} location, taking the common value where
the two views agree and \textsf{undef} where they differ.  Two views may
therefore each hold one edge into a node and compose to a view holding both,
which is exactly what \textbf{OW} denies.
\texttt{composition\_does\_not\_preserve\_WellFormed} in
\texttt{rocq/Compose.v} is the witness: two well-formed views, heaps agreeing
wherever both are defined, no shared stack, no shared free list, the same lock,
root, readers and bounding threads --- every side condition the figure attaches
is met --- whose composition fails \textbf{OW}.

\paragraph{The repair.}  The counterexample says the operation is wrong, not
that framing is impossible, and \texttt{WellFormed\_comp} gives a composition
under which \textbf{WellFormed} is closed.  Two changes, and the second is the
one worth reporting.

First, the heap is \emph{shared} rather than split: composable views agree on
it.  That is the right reading for this structure --- every thread reads the
whole of it, and a separating heap would leave a reader unable to see the path
it is traversing --- and it disposes of the counterexample, since agreeing views
cannot between them hold two edges into one node unless one of them already
did.

That alone is not enough, and the five invariants it leaves broken have a common
shape.  \textbf{IFL} relates an observation to a free-list entry; \textbf{WULK}
and \textbf{FNR} relate two observations; \textbf{FR} and \textbf{WF} relate a
stack binding to an observation.  Each is a relation between two pieces of
\emph{per-thread} state about one location, and composition is precisely the
operation that can put those two pieces in different views.  The other fourteen
either relate the shared heap to per-thread state or are monotone in the
observations, and they survive untouched.  So the split is not arbitrary: an
invariant is closed under composition exactly when it does not constrain two
pieces of split state at once.

The second change follows from that reading.  A location's auxiliary state may
not be split --- observations and free-list entry live in one view --- and a
view's stack names only locations it owns.  With those, closure holds for all
nineteen.  They are stronger conditions than the published figure attaches, and
they are what it is missing.

Finally \texttt{repaired\_composition\_is\_not\_vacuous}, because a composition
defined on nothing would make the closure theorem true and useless --- the
failure mode \textbf{RITR} exhibits, one floor up.

One detail confirms the reading.  The report's own proofs of the atomic actions
state their composition hypotheses as $\textsf{WellFormed}(\sigma, O_1, U_1,
T_1, F_1)$ and $\textsf{WellFormed}(\sigma, O_2, U_2, T_2, F_2)$ --- the
\emph{same} $\sigma$ for both views.  So the proofs already treat the machine
state as shared, and it is the figure that is out of step with them.  The repair
is not a design change; it is what the rest of the development was assuming.

\paragraph{Interference: a reading, not a defect.}  The same sentence asserts
closure under interference, and it is worth saying why that one is a different
kind of claim.  Composition is declared as a function $\bullet : \mathcal{M}
\rightarrow \mathcal{M} \rightarrow \mathcal{M}$, so landing in $\mathcal{M}$ is
an obligation and the counterexample above discharges it negatively.
Interference is declared as a relation $\mathcal{R} \subseteq \mathcal{M} \times
\mathcal{M}$, so it relates well-formed states to well-formed states by
construction, and the parenthetical restates the typing rather than adding to
it.

The formula in Figure \ref{fig:comp} does not on its own entail that, and a
reader mechanizing from the figure will get it wrong ---
\texttt{interference\_does\_not\_preserve\_WellFormed} shows a pair the formula
relates in which the successor is not well formed, because
$\mathcal{R}_0$ constrains the successor heap only by requiring that nodes
already observed as iterators or as the root stay allocated.  So $\mathcal{R}_0$
has to be read as intersected with $\mathcal{M} \times \mathcal{M}$.  That is a
reading, and we record it as one.

What the typing does \emph{not} settle is the direction that matters for
framing, which is whether $\mathcal{R}$ is permissive \emph{enough}: the
obligation for each action requires the frame to absorb that action's effect,
so every action's guarantee must be contained in the rely.  Nothing above
establishes that, and it is where the remaining risk sits.

\paragraph{Where framing should go instead.}  Both difficulties have one
source: every view carries a copy of the whole heap.  With a single writer that
is awkward; with several it is untenable, since every write must then be
absorbed by every frame and the rely degenerates to ``some other thread did
something that preserved the invariants''.  Repairing the operator and the
relation buys a single-writer proof whose structure does not survive the
generalisation.

\texttt{rocq/Triples.v} tests the alternative, which is the one
\texttt{IrisGhost.v} was already built for: put the heap in the invariant rather
than in the views.  Composition is then on the ghost state alone --- observations
and the free list --- which is genuinely separating, so the counterexample
cannot be written down; framing is Iris framing, with no bespoke monoid and no
hand-rolled interference relation; and the atomic-action lemmas become the
obligation to re-establish the invariant before closing it.

That last point is what the experiment checks, and it holds.  \textsc{Free}
restated against the invariant is \texttt{free\_update}, and its entire pure
content is \texttt{free\_preserves\_WellFormed} applied once: the action lemma
plugs in unchanged.  One thing was missing from the ghost layer and had to be
added, \texttt{fl\_delete}, which drops a free-list entry --- the operation
\textsc{Free} needs and which exclusive control of the entry is what licenses.
Framing is not mentioned in the proof because it is not needed there.

\textsc{T-UnlinkH} is the second, and the one that matters, because it changes
the observations as well as the heap and because what it changes is \emph{one
thread's} view of one location: the writer's iterator on the node being
unlinked, while a reader's iterator on that same node must survive untouched.
\texttt{unlink\_update} discharges it, and again the pure content is the action
lemma applied once.  The per-thread ghost state earns its place here rather than
in a design note --- the revocation is an update at the key $(z, l)$, and a
reader's observation at $(z, t)$ is simply a different key.

It needed two things the \textsc{Free} case did not, and both are worth
recording.  The invariant has to carry \textbf{ObsWF}, that an entry at $(o,t)$
records only observations of $t$ or the anonymous \rt{}; without it, revoking
the writer's entry would not rule out a writer-tagged iterator sitting in some
other thread's entry.  And the writer must hold \emph{exactly} its iterator
observation on the node --- the singleton fragment, not merely some fragment ---
which is what makes the revocation legal and what discharges the requirement
that every other observation survives.

Ten of the fifteen went through in that form: \textsc{Free},
\textsc{T-UnlinkH}, \textsc{T-WriteFH}, \textsc{ReadBegin}, \textsc{T-Alloc},
a reader's read, the three binding rules, \textsc{T-Insert},
\textsc{T-LinkF-Null} and \textsc{T-Replace}.  In every one the pure content is
the action lemma applied once and the ghost work is an update at a single key
--- two, for \textsc{T-Replace}, which promotes and demotes in the same step.

The other five change the ghost state at many keys at once.  \textsc{SyncStart}
gives every detached node a free-list entry; \textsc{SyncStop} turns every
\unlk{} observation into a \frbl{} one; \textsc{ReadEnd} and \textsc{WriteEnd}
retire a thread's entries; \textsc{WriteBegin} grants the incoming writer an
iterator on every reachable node.  The ghost layer had only single-key
operations, so none of them could be stated at all.

What they need turns out to be neither a new camera nor a new frame-preserving
update.  A map-wide change is an iterated single-key change, and what makes it
legal is owning every fragment it touches --- or the key being free, which is
the other way to earn the right to write there.  \texttt{tobs\_upd\_list} and
\texttt{fl\_upd\_list} are that induction, and it is four lines each.  Worth
noticing is what they do \emph{not} ask: that the intermediate maps satisfy
\textsf{WellFormed}.  The fold happens inside one opening of the invariant, so the
half-finished states never appear in it and only the endpoint must be well
formed.  That is what lets \textsc{SyncStop} recolour every observation in a
single step.

With those, all fifteen are in that form.  Three of the five are worth a
sentence each.

\textsc{SyncStop} reads like the one action that is not thread-local: every
\unlk{} observation becomes \frbl{}, which sounds like a change to every
thread's entries, and would mean a reader had to surrender its fragments for
the writer to recolour them.  It is not.  \textbf{WUNLK} says only the writer
holds a detaching observation and \textbf{ObsWF} says an entry records only its
own thread's, so the keys whose contents actually change are exactly the
writer's; everywhere else the recolouring is the identity and no update is
needed.  \textbf{WUNLK} is one of the two invariants this development added,
and this is where it pays for itself: without it the step is global.

\textsc{ReadEnd} and \textsc{WriteEnd} retire a thread's observations, and
``retire'' is deliberately not ``delete''.  The anonymous root observation may
be filed under any thread --- \textbf{ObsWF} permits it, and the counterexample
state \texttt{ce\_Og} is one where it is filed under a reader --- and the action
lemmas require it to survive a departure.  So the retirement keeps \rt{} and
drops the rest, which is what makes the survival condition provable instead of
an extra assumption about where the root lives.

\textsc{SyncStart} needed something the nineteen do not supply, which is
\S\ref{sec:defect-fld}.

\paragraph{A defect in the bridge, found the same way.}
Restating \textsc{ReadEnd} and \textsc{WriteEnd} turned up something worse than
a missing operation.  Both require \emph{every} variable of the departing
thread to leave scope.  The bridge represented scope as a finite set of
variable/thread pairs that are out of scope, and with variables drawn from an
infinite supply no finite set contains every variable of a thread.  The
hypothesis was unsatisfiable, so the two action theorems that carried it ---
proved, checked, axiom-free --- established nothing whatever.

\texttt{no\_finite\_scope\_set} is the witness.  Flipping the polarity does not
help: the initial state has nothing out of scope, which under the flip needs the
set to be everything.  Scope is now a predicate rather than a set, which costs
nothing --- it carries no ghost authority, the invariant merely quantifies over
it --- and makes both theorems mean what they say.

This is a defect in the mechanization rather than in the type system, and it is
reported here for that reason.  A vacuous theorem is the one failure mode a
proof assistant does not catch: it type-checks, it closes under the global
context, and it reports no axioms.  Only using the theorem finds it.

\paragraph{Choosing \texttt{phys}, and a defect in the state model.}
That leaves the physical state, which was a parameter and the physical step a
hypothesis.  Choosing \texttt{phys} is what turns a step on the invariant's
contents into a Hoare triple, and attempting it is what found the last defect.

The standard choice is a heap of points-to assertions, so that a thread's right
to write a field is a resource.  It was not available.  Allocation as published
gives the new node \emph{every} field, and field names are drawn from an
infinite supply, so one allocation puts infinitely many cells in the heap.  No
finite map represents such a heap, so there is no \texttt{gen\_heap} over it and
no camera of the usual shape either --- they are all built on finite maps.  And
without points-to a rule whose premise mentions the heap cannot be stated
outside the invariant at all, because the heap is bound inside it.
\texttt{no\_finite\_heap\_map} is the proof, stated about the published
definition, which the development retains under the name \texttt{alloc\_all} so
that the defect can be written down.

The repair is a class declaration: allocation takes the node's field list and
initialises exactly those fields.  The heap is then finite if it starts finite,
which \texttt{Represents\_upd}, \texttt{Represents\_free} and
\texttt{Represents\_alloc} establish for the three heap-changing operations.
The cost is small and local --- \texttt{alloc} gains an argument, heap-domain
closure is unchanged, and the allocation case gains one hypothesis, that the
field list holds every RCU field, which is what makes it a declaration rather
than an arbitrary subset and what the fresh node's denotation needs.

\paragraph{Ten closed triples: the whole write side, and the read side's two.}
With a finite heap there are points-to assertions, and with them eight actions
close completely: \textsc{Free}, all five heap mutations, \textsc{T-Alloc}, and
the binding rules --- which is to say the entire write side.  \textsc{ReadBegin}
and \textsc{ReadEnd} join them once the free list stops being anybody's
exclusive property, which is \S\ref{sec:inv-readside}.  Each is stated with no
parameters and nothing assumed about shared state: the invariant is opened,
every premise is \emph{derived}, the step is taken, the invariant is closed.

There are four kinds of premise, and each is discharged differently.
\emph{Observational} ones come from the thread's own ghost fragment ---
exclusive control of its entry at a node, and the agreement lemma turns that
into ``the node is observed \frbl{}''.  The typing side condition becomes a
resource.  \emph{Physical} ones come from points-to, the lock token and the
stack fragment, none of which existed before the field list made the heap
finite.  \emph{Path} premises come from a fold of lookups in the cell map the
thread owns; what that derivation does \emph{not} use is the point ---
no invariant, no uniqueness, no shape condition, because reachability by a
\emph{named} path is a local fact about the cells on it.

\emph{Negative} premises --- that something is absent --- cannot come from a
fragment at all, since a fragment says what is and never what is not.  They come
from invariants.  ``No free-list entry here'' is \textbf{FLD} with
\textbf{WULK}, or with \textbf{FNR} for a fresh node.  ``This fresh node is
unreachable'' is \textbf{FR} with the pinned root.  ``No other thread observes
this node as fresh'' is the repaired \textbf{WFresh}.  ``Every detaching
observation is the writer's'' is \textbf{WUNLK}.  ``This node is not detached''
is \textbf{WULK} and \textbf{FNR} together.  That is a concrete answer to what
the global invariants are \emph{for} in a separation-logic reading of this
system: they are exactly the facts a resource cannot carry --- and three of the
five just named are invariants this chapter adds.

\paragraph{Two premises of no such kind.}
The aliasing premise the unlinking rules carry --- no \frsh{} reference anywhere
points at the node being unlinked --- is negative and is not an invariant.  What
discharges it in the type system is the environment itself, whose denotation is
an intersection over all of it.  Here the invariant records that the fresh nodes
are \emph{enumerated}: every location observed \frsh{} lies in a set whose
authority the invariant holds.  The thread holds the matching exclusive fragment
and those nodes' cells, so the premise becomes
$\forall q,g \ldotp \mathit{val}(q,g) \neq z$ --- about values the thread owns,
mentioning no shared state, and exactly what the checker's
\texttt{noFreshPointsAt} decides by walking the environment.  The quantifier
does not disappear; it moves from the shared state, where nothing could
discharge it, to the thread's own resources, where the type system already had
it.  That is what translating a type environment into separation logic means.

\textsc{T-Alloc}'s six premises say the new location is unallocated,
unreferenced, pointed at by nothing, unobserved, unlisted and not the root.
Those look as though they need an operational semantics --- something must
\emph{produce} a fresh address --- and we expected to have to stop there.  They
do not.  An allocator is a function of the state, and the state is now finite in
every component, so the locations used for anything form a finite set and a
location outside it satisfies all six at once.  The field-list repair pays
twice: it made the heap representable, and it made the allocator constructible.

\paragraph{The type environment as a resource.}
The eight triples say the invariants survive and hand the thread its fragments
back.  Axiom soundness asks for more: that the post-state satisfies the
denotation of the post-type environment.  The pure half of that is already
proved --- there is a \texttt{post\_env} theorem for every write-side rule ---
so what was missing is the bridge, a reading of a type environment as a resource
from which the pure denotation follows.

The obvious reading, one resource per variable, does not work, and the reason is
worth stating because it is the same reason twice.  Two iterators may name the
same node, and two paths overlap in their prefix; a per-variable reading would
demand two exclusive fragments for one location and two owners for one cell.
The aliasing side conditions are exactly what rule that out, but they are
conditions on the \emph{environment}, not on the individual variable, so a
per-variable reading cannot see them.

What works is to own the resources \emph{once}, as maps, and make each
variable's requirement a pure condition on those maps.  The thread holds its
stack entries, its observations, its free-list entries and a set of cells; a
path is then a fold of lookups in the cell map, an alias is two variables
reading the same entry, and nothing is duplicated because a map has one entry
per key.  \texttt{writer\_env} is the resulting theorem: what the writer holds
implies the environment denotation, for every type the write side produces.  And
\texttt{EnvOK\_alloc} is the same reading across a step --- the environment
survives an allocation because the new node is one nothing else mentions, which
is the allocator's condition used a second time.

This has a consequence worth recording.  \emph{No fractions are needed on the
write side}: overlapping paths share cells by being the same map entry, not by
splitting ownership.  The read side still needs them, for the different reason
that a reader and a writer must hold the same cell at once.

And then one rule stated that way end to end.  \texttt{alloc\_typed} is
\textsc{T-Alloc} with a type environment on both sides: the thread hands in the
resources its environment reads as $\Gamma$ and gets back the resources its
environment reads as $\Gamma$ with the allocated variable rebound to
\texttt{rcuFresh}.  Nothing new is proved in it --- the step is the closed
triple's, the environment's survival is the monotonicity lemma, and the
freshness conditions the latter needs are the allocator's, read off the thread's
own maps because they are submaps of the invariant's.  It is the shape all
fifteen should have, and the first of them to have it.

The other seven rules do not grow the maps at a new location; they overwrite an
entry.  So what they need is not monotonicity but \emph{framing}: the
environment survives a step whose footprint it does not read.  And there are
only three kinds of footprint --- a cell, a location's observations, a stack
slot --- so there are three lemmas rather than seven.

The cell case is the one with content, and it is where the aliasing side
conditions land.  A path is a fold of lookups, so the cells it reads are a
concrete list, and a write outside that list cannot change the fold.  That is
\texttt{EnvOK\_cell}, and its hypothesis --- the written cell is on no path the
environment reads --- is exactly the question \texttt{MayAlias} decides in
\S\ref{sec:pathfrag}.  The decision procedure and the framing lemma are the
same condition, one decided and one proved, which is the tightest connection
between the two halves of this work that we have.

What is left for each rule is the variable it \emph{touches}, and that turned
out to need one thing the reading did not yet have: field maps.  Restricting
types to empty field maps made three of the post-types inexpressible ---
\textsc{T-WriteFH} produces a fresh node with a field written, and
\textsc{T-Insert} and \textsc{T-Replace} take field maps as premises --- so the
reading now carries them.  A field map is read on the thread's own maps too:
$N(f) = y$ is a lookup in the cell map whose target is a variable the
environment also names.  That is \texttt{FieldOK}, and it is what makes the framing
lemmas' hypotheses complete: a step must not touch a tracked field either.

With that, the touched variable is not a lemma but the definition supplied.  The
three cases that occur --- demotion to \unlk{}, to \frbl{}, and promotion of a
fresh node to an iterator one field beyond its parent --- are written out in the
development so the claim is checkable rather than asserted.  Only the promotion
has any content, and the content is that the new path is the old one extended:
the fold reaches the parent and then takes the cell the rule just wrote.

Making the rules compose with that reading removed the last two ad-hoc
resources.  The fresh nodes had been a separate assertion and the paths had been
one assertion per iterator; both are now conditions on the single cell map the
writer owns.  The path case is the instructive one: one assertion per iterator
cannot be right, because two iterators' paths overlap in their prefix and a
points-to is exclusive.  All five mutation rules are now stated with their paths
and their fresh nodes as pure facts about that map, which is exactly the form
the environment reading uses --- so the two sides of a rule are written in the
same language.

With that, a mutation goes end to end.  \texttt{link\_null\_typed} is
\textsc{T-LinkF-Null} with a type environment on both sides: the variables the
step does not touch are framed across it by the two lemmas, applied in the order
the step changes things so that both hypotheses are about the cell map the
caller started with; and the one it does touch gets its new type, which is
\texttt{TyOK\_promoted}.  The only work in the touched case is that the fresh
node's new path is the parent's extended by the field just written --- and that
its chain of iterator observations is the parent's chain plus one step, for
which a prefix of $\rho.f$ is either a prefix of $\rho$ or the whole of it.

None of that is about RCU.  It is the bookkeeping that connects a rule's
resources to a rule's types, it is the same bookkeeping for all five mutations,
and having done it once the remaining four are a transcription.

\textsc{T-UnlinkH} follows, and it is the transcription the first one was
supposed to make possible.  It touches three variables rather than one --- the
parent's field map is retargeted, the victim is demoted to \unlk{}, and the node
below is promoted to the parent's field --- and each is one of the three
toolkit lemmas, with the rest of the environment framed by the same two in the
same order.  It took no new ideas, which is what the toolkit was for.  Two
restrictions are worth naming: the parent and the promoted node are taken with
empty field maps, which is the case both worked examples use, and the victim is
assumed off the parent's path, which is uniqueness of paths and which the rule's
premises supply.

\textsc{T-WriteFH}, \textsc{T-Insert} and \textsc{T-Replace} follow without
incident, and with them every heap mutation and the allocation are stated with
the environment denotation as their conclusion.  We should be exact about the
other side, because the six are not uniform there.  \textsc{T-Alloc} and
\textsc{T-LinkF-Null} take the whole of $\Gamma$ as a hypothesis; the other four
take the untouched part of it together with the components --- the stack entry,
the observation, the path --- that $\Gamma$ supplies for each touched variable.
Those components are exactly what the denotation of the touched variable's type
unfolds to, so nothing is assumed that $\Gamma$ does not give; but the
hypothesis is $\Gamma$ unpacked rather than $\Gamma$, and packing it is work we
have not done.

The pattern is otherwise uniform: the untouched variables are framed by the two
lemmas, and the touched ones are supplied, each being one of three forms.  \textsc{T-WriteFH} is the smallest, changing no observation at
all, so only the cell framing applies and its one variable stays \frsh{} with a
field added.  \textsc{T-Insert} touches three --- the parent retargeted, the
fresh node promoted, the displaced node moved one step further down, which is
the promotion lemma applied to a cell the fresh node already held.
\textsc{T-Replace} is the two other rules at once: a promotion and a demotion,
with two observation entries changing, so the observation framing applies twice.

That the last three needed no new machinery is the substance of the claim.  What
varies between the rules is which of the three touched-variable forms occur and
how many times; what does not vary is anything about RCU.

Three more rules are stated that way, and two of them needed framing lemmas the
five mutations did not, for reasons worth separating.

\textsc{T-Free} is the only action that makes the thread \emph{give resources
up} rather than write them, so it is the only place the cell map has to come
apart --- into the reclaimed node's cells, which the step consumes, and the
rest, which comes back with the environment.  That also makes the path condition
the harder direction of the two: a write outside a path cannot change the fold,
but a \emph{deletion} outside it might make the fold fail rather than return
something else, so the invariance has to be proved separately
(\texttt{hstarC\_stable\_free}).  Its environment condition is
\texttt{AvoidsNode}, and the post-environment drops the freed variable rather
than retyping it to \undf{} --- not an evasion, because \undf{}'s denotation
asks for a fact about a node the thread no longer references, so its honest
resource reading is no resource at all, and what \undf{} \emph{means} is
\texttt{free\_post\_env}, where it belongs.

The binding rules needed the mirror image.  \texttt{EnvOK\_obs} frames an
observation write at a node the environment does not read, and a bind writes at
the node it does; but every condition a type puts on an entry is a membership,
so an entry that only grows carries all of them (\texttt{EnvOK\_obs\_grow}).
With that, \textsc{T-ReadS} and \textsc{T-Root} are one application each.

The reader's read is the ninth, and it is the one that is not a triple.  The
reader's types are read on a smaller condition than the writer's --- no paths,
no cells, which is \texttt{REnvOK} --- and that condition survives the read for
the same reason the binding rules work: the stack write is at the rebound
variable's slot and the observation write only grows the entry at the node
reached, so nothing has to be assumed about aliasing between the environment and
the node read.  What remains a hypothesis is the edge itself, for the
fractional-heap reason in \S\ref{sec:inv-readside}.

The last six are the protocol, and their environments are the smallest in the
file.  \textsc{ReadBegin} and \textsc{ReadEnd} have environments ToRCURead
fixes rather than derives --- the section starts with nothing and ends with
nothing --- so what is worth proving is that the two \emph{fit}: entering, the
reader gets a registration whose domain is the domain of its empty observation
map, which is exactly the shape the read consumes; leaving, it hands back
exactly the entries that domain names.  \texttt{read\_section\_typed} is the
two composed, which makes re-entry a matter of resources rather than of
argument.

\textsc{SyncStart} touches no observation, so the environment is unchanged and
all that is needed is that a free list keeping the entries it had carries it
(\texttt{EnvOK\_fl}); the free list is read by exactly one type.
\textsc{SyncStop} touches no free-list entry, so the environment changes
instead: every variable typed \unlk{} becomes \frbl{} (\texttt{syncenv}), and
every other type asks for an observation the recolouring fixes.  The certificate
the retyped variables need is the grace period's own rather than the
environment's, which is why it is a hypothesis of \texttt{EnvOK\_syncstop} and
discharged where the watermark is.

\textsc{WriteBegin}'s environment half is one lookup, and it is one lookup
because of the root observation.  The only type the writer starts with is
\rt{}, and \texttt{RootOK} is now a stack binding and nothing else; under the
old encoding it would have asked for the entry at the root to hold \rt{}, which
is precisely the entry \textsc{WriteBegin}'s bulk update overwrites.
\textsc{WriteEnd}'s is empty, for the same reason \textsc{ReadEnd}'s is, and it
is now a statement about resources rather than a convention because
\textsc{WriteEnd} deletes its keys rather than blanking them.

\paragraph{Chaining the rules, and what that found.}
Each typed rule was proved in isolation, and nothing had yet forced one rule's
postcondition to \emph{be} the next one's precondition.  Interface mismatches
are the class of defect that only appears under composition, so
\texttt{rocq/BST.v} runs the write side of the paper's binary search tree delete
--- the allocation and the splice --- through the rules as a chain, over a
concrete four-node tree.

It found one immediately.  \textsc{T-Alloc} derives four facts about the
location it invents --- that it is not the root, and that none of the caller's
three maps already mentions it --- and exported none of them.  In isolation that
is invisible, because the lemma's own conclusion does not need them.  Under
composition it is fatal: the environment comes back stated over the extended
maps, and without those facts \emph{no lookup in them can be discharged}, so the
next rule cannot start.  Exporting them is the repair, and it is the reason a
rule that invents a name is different in kind from one that does not.

Nothing similar is expected of the other five, and for a reason worth stating:
they invent nothing.  Their output maps are the caller's own, extended at keys
the caller supplied, so the caller already knows everything needed to look
things up in them.  Allocation was the only rule that could have this defect.

The chain then found something sharper, and it is not about the mechanization.
Three of the four steps compose --- the allocation and the two writes that build
the replacement, in \texttt{bst\_build\_replacement}.  The fourth,
\textsc{T-Replace}, cannot join them, and not because of any interface
mismatch: its hypotheses are \emph{inconsistent} with the allocation's.
\textsc{T-Replace} is stated with ``every field is an RCU field'', the
simplification the mutation rules inherit from the report;
\textsc{T-Alloc} needs the class declaration the finite heap forced, that every
RCU field is one of the node's declared fields.  Together those say every field
name is declared, and there are infinitely many field names and two declared
ones.  \texttt{class\_excludes\_all\_rcu} is the two-line proof.

So the two rules cannot both apply in one program, and the delete is exactly a
program that needs both: it allocates a replacement and then splices it in.  We
had reported this incompatibility as a caveat alongside the field-list repair.
It is better as an obstruction: \emph{the paper's own worked example cannot be
typed} as things stand.

Tracing where the all-RCU assumption is actually used then said what the repair
is, and it is not the one we first guessed.  The mutation proofs reach for it in
six places, and \emph{every one of them applies it to an edge already in hand}
--- they are proving something about a node that some field points at.  So what
they need is not that every field is an RCU field, but that every field
\emph{holding a reference} is: the heap's reference structure lives in RCU
fields, and a scalar field holds a scalar.

That is a weakening, so every proof that used the old assumption still goes
through, and it is consistent with a class declaration --- which is what the
worked example needs.  We have made it.  Two rules now say something they should
have said all along: \textsc{T-WriteFH} and \textsc{T-LinkF-Null} write a field
that did not previously hold a reference, so they must state that the field they
write is an RCU field; the other three inherit it from the edge they are
replacing.

With the repair in place the splice is no longer blocked by the rules'
premises --- and it is blocked by something else, which the chain found in turn
and which is in this development rather than in the rules.  Replacing
\texttt{current} by the fresh node re-routes every path that ran through
\texttt{current}: \texttt{currentL} is still at $l.l$ afterwards, because the
fresh node mirrors what it replaces, but the \emph{intermediate} node on that
path is now the fresh one.  \texttt{rcuItr}'s chain-of-iterators condition is
about those intermediate nodes, so it has to be re-derived rather than framed.

The framing lemmas do not do that.  \texttt{EnvOK\_obs} frames a variable across
a change to one location's observations provided no path in the environment
passes through that location, and here two of them do.  The condition is not
merely unproved: \texttt{replace\_cannot\_frame} shows it is false.  So the
splice needs a framing lemma that knows about mirroring --- the paths are
unchanged as paths, and the nodes they pass through are swapped one for one,
which is what \textbf{Mirrors} is for in the pure layer.  We have now built its
counterpart in the environment reading.

\texttt{hstarC\_mirror} is the content: a fold that reached a node before
reaches its swap afterwards, because at the one cell the rule rewrote the target
moves from the old node to the new, and everywhere else the mirror makes the two
indistinguishable.  It needs one thing the other framing lemmas do not --- that
the replaced node has a single predecessor, which is \textbf{OW} --- and that is
the first place \textbf{OW} is used to frame rather than to establish an
invariant.  \texttt{EnvOK\_mirror} lifts it to the environment: a variable below
the replacement keeps its type, its path is the same path, and the intermediate
node it passes through is the new one, which the promotion has just made an
iterator.

So the framing story has three shapes, not two.  A step may leave what the
environment reads alone, which the first two lemmas cover; or it may move what
the environment reads to somewhere indistinguishable, which is this one.  The
second shape is what makes read-copy-update work at all, so it is fitting that
the framing lemma for it is the one that needs No-Sharing.

\textsc{T-Replace} is stated against that lemma now rather than against the
other two, so the splice is no longer blocked at either end: not by the rules'
premises, which the reference repair fixed, and not by the framing, which the
mirror lemma supplies.  With both repairs in place the four steps compose:
\texttt{bst\_delete\_two\_child} runs the allocation, the two writes that build
the replacement, and the splice as a single chain over the concrete tree,
starting from the entry environment and ending with \texttt{current} typed
\texttt{unlinked} and the fresh node in its place.  The write half of the
paper's worked example is mechanized.

\texttt{rocq/LinkedList.v} does the paper's other example, and the reason to do
it is different.  The delete is a single unlinking, \texttt{prev.Next =
current.Next}, so it finds nothing the tree did not; what it checks is that the
machinery is about the rules rather than about trees.  The two files differ in
what one would want them to differ in --- the tree needs four steps and the
mirror framing lemma, the list needs one step and the ordinary one --- while the
rule applied is the same rule at the same interface, with no list-specific
machinery on either side.  That is the claim \S\ref{sec:relatedwork} makes about
reuse across singly-linked and tree structures, checked rather than asserted.
The surviving variable also differs in the way that matters: in the list it sits
\emph{above} the deleted node, so ordinary framing suffices, where in the tree
it sits below the replaced one, which is what forced the mirror.

Two things about the tree's instantiation are worth reporting, because they are
what the general lemmas could not say.  The mirror condition is not an abstract
hypothesis there: it is discharged by enumeration, and the clause that does the
work is the one saying no \emph{prefix} of a surviving variable's path reaches
the fresh node --- three folds per variable in a four-node tree.  And the
single-predecessor fact, which is \textbf{OW} in general, is a finite check on
the tree's own cells, which is the sense in which the worked example pays for
the invariant rather than assuming it.

We tried two larger repairs first and record why they were wrong, because the
wrong ones are instructive.  Guarding \textbf{ULKR} by RCU-ness, so that it
matches \textbf{OW}, makes the invariants agree --- but \textbf{HD} needs the
same treatment for the same reason, and the reachability form of \textbf{ULKR}
chains along a path, so the guard becomes ``every field on the path is an RCU
field'' and propagates into every use.  Both of those are repairs to what the
invariants \emph{say}.  The defect is not there.  It is in the assumption the
proofs make to connect them, and weakening that assumption leaves every
invariant as published.  That is the second of the two rule defects: not that
\textbf{ULKR} and \textbf{OW} disagree, but that the published proofs reconcile
them by assuming a system with no scalar fields at all.

Two things this reading turned up.  The denotation of \texttt{freeable} was
stated as $F(o) = \lambda t. \bot$ --- an equality of \emph{functions} --- and
in the ghost encoding the entry is $\lambda t.\, t \in S$ for a finite set $S$;
proving those equal needs functional extensionality, which this development does
not assume.  A denotation that cannot be established without an axiom is a
defect in the denotation, and it is stated extensionally now.  And
\texttt{undef} has no resource reading at all: its second conjunct is about the
free list at a node the thread no longer references, which is state it does not
own.  That conjunct belongs in an invariant rather than in a type, and we report
it rather than move it.

\paragraph{The read side, and what is left.}
\label{sec:inv-readside}
The read side and the protocol remain, and the two obstacles are particular
rather than general.

We said in an earlier draft that the reader's read could not be stated at all,
because the points-to assertion is exclusive and a reader must read cells the
writer also holds; and that the repair was a fractional heap.  That was wrong,
and correcting it removes an item rather than adding one.

A reader's iterator is not a claim about the heap.  \texttt{D\_rcuItrR} is a
stack binding, an \itr{} observation of the reader's own, and the bounding
conjunct --- no path, no cell.  So the reader's environment follows from
fragments of the stack and the observation map and nothing else, which is what
\texttt{reader\_env} says: there is no heap authority anywhere in its statement,
and none in \texttt{reader\_acquire\_update} either.  What the read \emph{does}
with the heap is consult the authority inside the invariant to find the
successor and carry nothing out, and that needs no fragment, so it needs no
fraction.

What the read does need is \textbf{IFL}'s obligation: a reader acquiring a
reference to a node already on the free list must be one of the threads its
grace period waits for.  That is a real obligation and it is not discharged by
owning anything --- and working out what discharges it is what turned up
\S\ref{sec:defect-samesnap}.  With the entries agreeing it is
\texttt{read\_bound}, four lines from \textbf{FLR} and \textbf{IFL}.  The read
needs one thing more, that the node it reaches is not \frsh{}, and that is
\S\ref{sec:defect-frw}; with \textbf{FRW} it is \texttt{read\_not\_fresh}, one
line.

Both conjuncts are now in the Iris invariant rather than hypotheses, and with
them the reader's read is a step on the invariant's contents,
\texttt{read\_update}.  \textbf{SameSnap} did not need the mechanical pass the
other added conjuncts took: it is a consequence of the epoch representation,
since the free-list entries are derived from the stamps and every stamp is at
one epoch (\texttt{HSS\_SameSnap}).  \textbf{FRW} did need the pass, because a
field write creates an edge and so has something to prove at each of the five
mutations --- in each case that the edge's new target is not fresh, which the
two linking rules get from promoting it in the same step and the others from
\textbf{FNR}.

What \texttt{read\_update} establishes is everything observational about the
read.  The bounding obligation is discharged, as above.  So is the one that
looks smaller and was not: the reader has to grant itself an observation at a
node it has not observed, and the ghost update for that has to know whether it
already has an entry there.  Nothing decided that while the root observation
lived in a thread's entry --- ``this thread observes nothing here'' did not
imply ``this thread has no entry here'' --- and it is decided now by the
reader's own registration cell.  That is the second half of removing
\textbf{RTO}, and it is why that removal is reported as a repair rather than as
tidying.

What is left between \texttt{read\_update} and a Hoare triple for $y := x.f$ is
one thing, and it is not observational: a reader must \emph{learn} what the
field holds, and the points-to assertion is exclusive.  The correction above
stands --- the reader's environment needs no fraction, and neither does the
step --- but the cell read does, and that is a fractional or snapshot heap.  It
is the only obstacle left on the read side, and it is standard.

\textsc{WriteBegin} asks that the lock is free, and \textsc{SyncStop} that no
thread is still bounding.  Those are the two points where the protocol
\emph{tests} or \emph{waits on} shared state rather than acting on state it
owns, and neither is a fact any thread can hold.  They are what an operational
semantics is for --- a compare-and-set and a wait loop --- and they are the only
two places in the system that need one.

The state the other four rules need is now built, and building it is what found
\textbf{BR} (\S\ref{sec:defect-br}) and \textbf{WUNLKW}
(\S\ref{sec:defect-wunlkw}).  It took three attempts, and what each one failed
at is more informative than the answer, because each failure is a premise the
camera could not express.

The first made a reader's state a \emph{set} token --- the disjoint-union
camera, one element per active reader.  It certifies membership and nothing
else, and \textsc{ReadBegin}'s premise is that the thread is \emph{not} already
a reader.  A protocol whose rules have negative premises needs the state, not
the membership.  The second made it a per-thread two-valued cell, owned by the
thread whether it is inside a critical section or outside one, so that both
directions follow from what the thread holds.  That is enough for
\textsc{ReadBegin} and not for \textsc{ReadEnd}, for the reason the next
subsection is about: a boolean cannot distinguish a reader's new critical
section from its old one.  The third, which is what the development has, is a
cell holding \emph{when} the thread registered and \emph{which nodes} it
observes --- an epoch and a finite set.  The epoch is what re-entrancy needs and
the set is what \textsc{ReadEnd}'s enumeration of its own observations needs.

$R$ becomes a finite map in the process, for the same reason the heap and the
stack did; the thread set was already finite, so this is the model catching up
with itself.  The bounding set needs no camera of its own: it is pinned by
\textbf{BR} and derived from the registrations, and \textsc{SyncStart} is the
step that establishes it.

\textsc{ReadBegin} is now a closed triple, \texttt{read\_begin\_atomic}, and it
is the first on the read side.  What makes it closed is worth stating, because
it is what the two new invariants were for.  The action lemma
\texttt{read\_begin\_update} carries both of \textsc{ReadBegin}'s premises as
hypotheses --- that the entering thread is not the writer, and that it holds no
detaching observation.  The triple proves them instead: the first from the
thread's own cell, since the invariant now says the lock holder has no cell at
all; the second from \textbf{WFreshW} and \textbf{WUNLKW} together with the
first.  A rule whose premises are all discharged rather than assumed is the
difference between mechanizing a lemma and mechanizing a rule.

\textsc{ReadEnd} is the next rule, and it did not close for a long time.  The
reason is worth reporting in full, because the repair changed the model rather
than the proof.  It is the single step in the system that writes state another
thread owns: the free-list snapshots are the writer's, created by
\textsc{SyncStart}, and a departing reader rewrites every one it appears in.
The resource-level lemma has to hold all of them, which is the wrong resource
for a per-thread act.

The obvious economy is to leave the snapshots alone and read them against the
threads still running --- a thread that has left its read section bounds
nothing, so intersecting a snapshot with $R$ ought to do the work of removing
the thread from it, and \textsc{ReadEnd} would then touch nothing but its own
reader cell.  We tried that repair and it is unsound, because of re-entrancy: a
thread that leaves its read section and enters a new one is in $R$ again, and
the intersected reading puts it back into every snapshot it was ever in,
including grace periods that began after it left.
\texttt{reentry\_breaks\_the\_intersected\_reading} is the witness.  What makes
it worth having tried is where the failure appears --- not at \textsc{ReadEnd},
which the change was about, but at \textsc{ReadBegin}, which the change did not
touch.

The repair the development already recorded is to make a snapshot's value a
disjoint union of tickets, one per bounding thread, so that a reader owns its
own membership and spends it.  That works, and re-entry cannot restore a spent
ticket.  But there is a better one, and \texttt{rocq/Epochs.v} works it out:
stop naming threads in snapshots at all.

Let a registration carry \emph{when} the reader registered, and a snapshot carry
\emph{when} the node was unlinked.  A snapshot's members are then the readers
that registered no later than its grace period began --- derived, not stored.
This is what an implementation's grace-period counter does, and what the
comparison with Tassarotti's QSBR verification describes;
the question is what it is worth as a model.

Three things, and each is a theorem.  \textsc{ReadEnd} becomes a write to the
departing thread's own registration and nothing else
(\texttt{e\_end\_writes\_only\_its\_own}), and that single write produces all
three updates the published rule performs as separate ones --- the thread leaves
$R$, leaves $B$, and leaves every snapshot (\texttt{e\_end\_is\_read\_end}).
That last is an equation, not an analogy: the free list after the deletion is
literally \texttt{read\_end\_F} applied to the free list before it, so
\texttt{e\_end\_realises\_read\_end} says the transition
\texttt{read\_end\_preserves\_WellFormed} is stated over is exactly what
deleting a registration produces.  The epoch state is a representation of the
same system, not a different one with similar properties.

Entering joins nothing (\texttt{e\_begin\_joins\_nothing}): a reader registers
at the current epoch and every running grace period was stamped earlier, so it
is in no snapshot, and the free list is therefore untouched
(\texttt{e\_begin\_keeps\_the\_free\_list}) --- which is the property the
published model has to arrange by removing threads as they leave.  And
\textsc{SyncStart}'s entry is the current reader set
(\texttt{e\_sync\_start\_snapshot}), which is what the published rule says it
is, so all three steps correspond and the only difference is which of them has
to write what.

The bookkeeping also still satisfies the invariants the snapshot model is stated
over --- \textbf{RINFL} is \texttt{e\_RINFL}, and \textbf{BR}, which the
published set is missing entirely (\S\ref{sec:defect-br}), is immediate.

The ownership objection has a second end, and the epochs answer that too.
\textsc{Free} needs to know a snapshot is empty, and under the published model
it knows by \emph{owning} the entry --- which is what made the entry the
writer's resource and the reader's write illegal in the first place.  Here it
need own nothing, because quiescence past an epoch is \emph{stable}: once every
registration is later than $e$, every registration stays later than $e$, under
all four steps (\texttt{quiescence\_is\_stable}).  Registrations only get newer,
a reader entering at the current epoch and the counter never going backwards.  A
stable fact needs no exclusive ownership in the resource reading --- it may be
duplicated and carried --- so \textsc{SyncStop} hands \textsc{Free} a
certificate rather than a resource (\texttt{e\_free\_premise}), and the free
list stops needing to be anybody's property at all.

That is the whole of the objection answered: the reader writes only its own
registration, and the writer reads only a fact it already holds.

The re-entrancy failure is repaired, and the counters say exactly where the
repair lives: \texttt{reentry\_safe} turns on the epoch being incremented when
the grace period starts.  \texttt{increment\_is\_what\_makes\_reentry\_safe} is
the contrast --- the same three steps against a \textsc{SyncStart} that stamps
without incrementing, where the returning reader rejoins.  So the increment in
an implementation's counter is not an implementation detail: it is the thing
that distinguishes a reader's new critical section from its old one, and without
it the model has no way to tell them apart.

What it costs is a change to the state model, where the ticket costs only a
camera.  We took the epochs, and the development is now built on them.

That also discharges something \S\ref{sec:relatedwork} asserts of this paper.
It says a variation on our semantics would keep a vector of per-thread counts in
place of the reader set and a snapshot of counts in place of the bounding set,
and that ours is ``simpler than this alternative, while also equivalent''.  The
alternative is \texttt{rocq/Epochs.v}, and \texttt{Sim} is the relation: the two
states agree on who is reading, who is bounding, and what the free list is --
the published model's three components, recovered from the counters, which have
more structure than that.  \texttt{sim\_read\_begin},
\texttt{sim\_read\_end}, \texttt{sim\_sync\_start} and \texttt{sim\_sync\_stop}
are the four steps across it.

The one worth reading is \texttt{sim\_sync\_stop\_guard}.  The published rule
blocks until no thread is still bounding; the counters block until every reader
registered no earlier than the current epoch.  The theorem is that those are the
same condition --- which is the sentence the related work asserts, and the whole
of what ``more clearly equivalent to these implementations'' was claiming.

The ghost free list is stamps --- one epoch per node awaiting reclamation,
owned by the writer --- and the threads a snapshot contains are derived from the
registrations, owned one cell each by the readers.  The invariant binds the
logical free list and ties it to the two.  With that,
\texttt{read\_end\_atomic} closes.  The departing reader hands in its
registration and the observation entries its own cell enumerates, and gets its
registration back empty; its three premises are the enumeration having no
repeats, every entry in it being the thread's, and its covering the domain the
cell records, and the free list is not mentioned at all.  The invariant supplies
the converse --- that the domain covers every entry of the thread's carrying an
observation of its own --- and \textsc{ReadBegin} establishes it empty on entry,
from \textbf{WITR} for the iterator case and \textbf{WFreshW} and \textbf{WUNLKW}
for the three detaching ones.  The three invariants this work adds are what make
the read side's two rules closable.

One condition had to weaken on the way, and the reason is instructive.  The
retirement's coverage condition is now over entries carrying an observation
\emph{tagged with the departing thread}, not over all of the thread's entries.
It has to be: \textsc{ReadEnd} leaves the root observation behind --- it is
anonymous, and not the departing thread's to drop --- so a thread that has ended
a read section has entries it owns nothing about, and demanding that it
enumerate those would put the invariant's map straight back into its premises.

\textsc{Free}'s side changes with it.  A stamp is not enough to free a node;
what licenses the free is the certificate that the grace period past that stamp
is complete.  That threads through the environment reading as a pure parameter
beside the scope condition, and \frbl{} now means ``stamped at $e$, with a
certificate for $e$'' --- which is what the type meant all along, written where
a rule can use it.



That last point is not only ours.  The paper's own comparison with Tassarotti's
QSBR verification observes that implementations are structured so each reader
thread keeps its own \emph{fragment} of the shared state, updated only by that
thread and read by the write-side primitives.  That is the ticket, described as
an implementation technique.  The model abstracted it away, and \textsc{ReadEnd}
was the one place where the abstraction cost something: it turned a thread's
update to its own fragment into an update to a structure the writer owns.  The
model does not abstract it away any more.

The two testing points are mechanized as far as a development without a
semantics can take them, and further than we expected.  Both tests are over
finite data --- the lock is an option, the registrations are a finite map --- so
the test is decidable inside the proof and the attempt needs no oracle.
\texttt{sync\_stop\_attempt} is the wait's body: it decides whether every
registration still standing is later than the node's stamp, and either takes the
certificate or does not.  The certificate is a watermark fragment, which is
persistent because quiescence past an epoch is stable, and
\texttt{wm\_lb\_entry\_empty} is what it licenses --- the snapshot is empty,
which is the conjunct \frbl{} asks for.  The disjunction \emph{is} the loop
body.  Which branch is taken is the implementation's scan; that the left one is
eventually taken is the part that needs a semantics, and it is now the only
part.

\textsc{WriteBegin}'s test is decidable in the same way, and is not the obstacle
there.  What stops it being a triple of this shape is that it also takes an
\itr{} observation on every reachable node, which \textbf{UNQRT-b} is the reason
for; that is a bulk ghost update needing the reachable set enumerated, the same
premise \texttt{sync\_start\_update} carries.  We record that rather than state
a lemma that would claim more than it proves.

Termination reduces further than we expected, and to something that is not the
type system's to prove.  The threads \textsc{SyncStop} waits for are exactly the
snapshot at the node's stamp --- a finite set --- and each \textsc{ReadEnd} by a
member removes that member and nothing else (\texttt{wait\_shrinks}).  So the
wait ends after at most as many \textsc{ReadEnd}s as the snapshot has members,
and once they have all ended, in any order, the test passes
(\texttt{wait\_terminates}); \texttt{wait\_then\_free} carries that the last
step to what \frbl{} asks for.  What is left is fairness --- that a reader
inside a critical section eventually leaves it --- which is a property of the
client and not of the type system, and which no invariant of the shared state
could supply.  The lock's loop is the same shape and shorter: it ends when the
holder's critical section does.

Both conjuncts are carried in the invariant now.  \textbf{FRW} took the
mechanical pass and the five mutations are where that work is, a field write
being the only thing that creates an edge; \textbf{SameSnap} took none, because
the epoch representation gives it (\texttt{HSS\_SameSnap}).

The decision the reader's read wanted has been taken, and it went the way
\S\ref{sec:defect-obs} says.  The rule gives the reader an observation at a node
it may never have observed, and to grant one it must hold that node's entry or
know the entry absent.  Its cell records which nodes it observes, so the first
was always available for those; the second was not, because \textsc{ReadEnd}
left the \rt{} observation behind.  Reading \rt{} off the structure instead lets
\textsc{ReadEnd} delete its keys outright, makes the invariant's domain clause a
statement about entries rather than about tagged observations, and closes the
question: a node outside the cell's domain has no entry of the reader's.  That
is \texttt{read\_update}.

\paragraph{What the machinery is for.}
Everything above is machinery, and it is worth stating the property it exists
to establish, because once the invariants are right it is four lines and it uses
exactly two of them.

A variable typed \frbl{} names a node the writer is about to reclaim.  What
memory safety needs is that reclaiming it pulls the rug from under nobody.
\textbf{IFL} says a thread observing a node as an iterator is in that node's
free-list entry, and the \frbl{} denotation says the entry is empty, because the
grace period has completed; so nobody observes the node as an iterator, which is
\texttt{freeable\_is\_unobserved}.  \textbf{RWOW} says a live variable's
referent carries an observation --- an iterator one, or a detaching one
belonging to the lock holder.  The first is now impossible and the second
identifies the holder as the writer itself, so
\texttt{no\_live\_reference\_to\_a\_freeable\_node}: the only live reference to
a node about to be freed is the writer's own.
\texttt{free\_invalidates\_nobody} states it the way a client would ask it.

Two invariants, and both are ones the published set has --- which is the right
outcome.  The sixteen we add are not there to prove this; they are there to make
the fifteen actions preserve the set that proves it, which is the division of
labour a global invariant set is supposed to have.

Said over a state, that is weaker than memory safety asserts, which is about
runs.  The gap is one induction, and \texttt{reachable\_WellFormed} and
\texttt{safety} close it: the invariants are inductive along executions, so at
\emph{no point} in any run from a well-formed start does a thread other than the
writer hold a live reference to a node the writer is about to reclaim.

\texttt{lstep} is the relation, one constructor per action carrying that
action's hypotheses, so \texttt{lstep\_preserves} is fifteen applications and
nothing else; \texttt{system\_is\_safe} is the theorem above stated at it.  That
is the whole system rather than one action, which is what the parameterisation
was for.

\paragraph{How much of this is about RCU.}
The safety property above uses two of the twenty invariants, \textbf{IFL} and
\textbf{RWOW}, and neither of the two theorems mentions a grace period.  That
is worth following up, because \itr{} --- ``thread $t$ has announced it may
access $o$'' --- is what a hazard pointer \emph{is}, and if the argument does
not know how reclamation is decided then the invariants are not twenty facts
about RCU.

We first ask how much of the system touches reclamation at all, and the sharpest
way to ask it is mechanical: change the free list to \emph{anything} and see
which invariants survive.  Twenty-one of the twenty-six conjuncts do, and the
proof of each is \texttt{exact} --- they are not preserved by an argument, they
are unchanged, because the free list does not occur in them.  The five that do
not survive are exactly \textbf{IFL}, \textbf{FLR}, \textbf{RINFL},
\textbf{FLD} and \textbf{SameSnap}, and one witness state breaks all five.

Then we abstract the discipline.  It is a record with two components --- ``$t$
is recorded as a potential accessor of $o$'' and ``$o$ may be reclaimed'' ---
and two conditions: that an observation is recorded, and that the guard means
nothing is recorded.  The two safety theorems go through against that interface,
and they are the \emph{same proofs} with every mention of a free list gone.

RCU is one instance: the recording relation is ``$t$ is in $o$'s entry'', the
condition is \textbf{IFL} exactly, and the original theorem falls back out,
which is the check that the abstraction weakened nothing.  Hazard pointers are
the other: recording is set membership, so the guard condition holds by
construction, and the invariant a hazard-pointer scheme maintains by publishing
before it dereferences is \textbf{IFL}'s analogue with the indices swapped ---
\textbf{IFL} puts the \emph{thread} in the node's entry, this puts the
\emph{node} in the thread's set.

Two cautions, because the result is easy to overstate.  This is the safety
argument, not the whole system: the five free-list invariants do real work
elsewhere --- \textbf{FLR} and \textbf{SameSnap} discharge the reader's read,
\textbf{FLD} is \textsc{SyncStart}'s, \textbf{RINFL} ties the entries to the
bounding set --- and a hazard-pointer system would need analogues of those or
would do without the rules that use them.  And what carries over is the
argument, not the implementation: hazard pointers must publish before
dereferencing and re-validate after, which is a memory-ordering obligation this
development does not model, being sequentially consistent.  The publication
invariant is stated over the logical state, and is therefore exactly what a
weak-memory account would have to earn.

Progress deserves a word, because the obvious statement is not the right one.
``Some step is always possible'' is false and should be --- a state in which no
thread can move is a deadlock, and whether one is reachable is a question about
the program.  \texttt{lstep} relates states, not configurations, so
relation-level progress is not statable at that level and would not mean what it
sounds like.  What progress means here is per-thread and per-primitive: that a
well-typed thread about to execute a primitive finds its side condition true,
which is what the six guards are, and which \S\ref{sec:alglave} treats.

\paragraph{Against Alglave et al.'s requirements.}
\label{sec:alglave}
The prose above argues that this model meets the fundamental requirements for an
RCU implementation.  Two of the four are now theorems rather than arguments, and
it is worth saying which invariants they need.

\emph{Grace period and memory barrier.}  The essential clause of Alglave et
al.'s correctness criterion is that reader critical sections do not span grace
periods.  A section spans one if it began before the grace period began and was
still running when it completed; the grace period for a node stamped at $e$
completes exactly when every registration still standing began later than $e$,
so \texttt{no\_section\_spans\_a\_grace\_period} is immediate from that
condition and \texttt{reclamation\_waits\_for\_overlapping\_readers} states it
about the reclamation.  The content is not in those statements but in what
establishes the condition, which is \texttt{wait\_terminates}, and in what it
licenses, which is \texttt{wait\_then\_free}.

\emph{Publish-subscribe.}  A freshly allocated node cannot be observed by a
reader until it is published.  \texttt{readers\_cannot\_see\_unpublished} needs
two invariants, and one of them is the one the reader's read forced
(\S\ref{sec:defect-frw}): \textbf{FNR} gives that a fresh node carries no
iterator observation, so no reader observes it, and \textbf{FRW} gives that it
has no incoming edge, so no reader can walk to one.  Without the second the
guarantee is half of itself --- a reader could not be \emph{holding} an
unpublished node, but nothing said it could not reach one.  That a published
correctness criterion turns out to need an invariant the published invariant set
does not have is the sharpest form the mechanization's finding takes.

\emph{Unconditional execution.}  We first wrote that this one needed no proof,
on the grounds that the primitives are total functions on the machine state.
That was wrong, and reading the semantics of \S\ref{sec:semantics} again is what
shows it.  Five of the six carry a side condition: \textsc{ReadBegin} is written
with $\mathit{tid} \neq \ell$ and with $R \uplus \{\mathit{tid}\}$ on the right,
so it requires the thread to be neither the writer nor already a reader;
\textsc{ReadEnd} is written with $R \uplus \{\mathit{tid}\}$ on the \emph{left},
so it requires the thread to be a reader; \textsc{SyncStart} requires an empty
bounding set; \textsc{WriteBegin} requires the lock free; \textsc{SyncStop}
requires the bounding set empty.  Only \textsc{WriteEnd} is unguarded.

So there is something to prove, and it is the right thing: not that the guards
are trivial but that a \emph{well-typed} thread never finds one false.  For the
reader's two that is a theorem.  \texttt{read\_begin\_guard} and
\texttt{read\_end\_guard} discharge them from the thread's own registration
together with the invariant's clause that the lock holder holds none --- a
thread with an unregistered cell is neither the writer nor a reader, which is
\textsc{ReadBegin}'s condition, and a thread with a registered one is a reader
and still not the writer, which is \textsc{ReadEnd}'s.
\texttt{write\_end\_unconditional} is the one that really is nothing.

The writer's three are not discharged, and \texttt{writer\_guards\_are\_not\_%
invariants} is the proof that they cannot be: a state satisfying every invariant
can have the lock held and a grace period running, and then all three are false
at once.  Two of them are the blocking ones, and blocking rather than failing is
what the requirement asks.  \textsc{SyncStart}'s is discharged by the protocol
rather than by any state --- \textsf{Sync} is \textsc{SyncStart} followed by
\textsc{SyncStop}, so a writer that has completed its previous grace period
finds the bounding set empty.

\emph{Read-to-write upgrade.}  Not provided, as \S\ref{sec:relatedwork} says: a
performance optimisation, orthogonal to memory safety.

\paragraph{From one model to any implementation.}
Everything above is about \emph{one} abstract model.  A client compiled against
a real RCU is safe only if that RCU refines the model, and refinement proved
once over an interface is one theorem, where refinement argued per
implementation is one argument per implementation forever.
\texttt{rocq/Refine.v} is that interface: an implementation is a state space
with the four protocol actions on it and three observations --- who is reading,
who is bounding, what the free list says --- and the obligation
\texttt{Refines} is five clauses, one per action plus one saying the
implementation's grace-period guard decides the model's.

The interface was read off a proof rather than invented.  \texttt{Epochs.v}
already contained a second model of the same protocol, together with a relation
and one theorem per action across it; \texttt{ISim} is that relation with the
epoch state made abstract, and \texttt{epochs\_refine} assembles the first
instance from the five theorems that file already had, unchanged.  That the old
proof slots into the new interface without adjustment is the evidence that the
interface is the right shape rather than a shape.

The fifth clause is the one with content, and it is worth saying why.  Four of
the five are about single steps agreeing, and an implementation that got one
wrong would be wrong obviously.  An implementation can satisfy all four and
still be unsafe, because the model's \textsc{SyncStop} does not itself check
anything --- the checking is in when it is allowed to run.
\texttt{eager\_is\_not\_a\_refinement} is the epoch model with its guard
replaced by ``always'': it satisfies the four step clauses, they being literally
the same theorems, and fails only the guard.  The condition is the safety
property, not a tidiness clause attached to four equations.

Reclamation is the clause with content, and it is the one the protocol exists
for.  Two clauses cover it: one is bookkeeping, that freeing removes the
free-list entry; the other is Alglave et al.'s first requirement stated where it
bites, that an implementation may free a node only once the node's entry is
empty.  Nothing constrains \emph{how} an implementation decides that --- the
counter model compares counters, something else may count quiescent states ---
only \emph{when} the answer may be yes.  A fifth clause belongs with it:
\texttt{ref\_snapshot}, that the grace period waits for exactly the threads now
reading.  \texttt{ref\_guard} says when the wait may end; this says what it was
waiting for, and the two together are requirement~1 at this interface.

The transfer then follows.  \texttt{xstep} pairs an implementation state with
the observation bookkeeping, carrying exactly the hypotheses the step relation
of \S\ref{sec:inv-mech} already had plus the implementation's guard, and the
two halves come from different places: the implementation supplies the protocol
half and the typing derivation supplies the observation half.
\texttt{xrun\_WellFormed} runs the pair, and \texttt{xrun\_safe} is the
payoff --- a run of a conforming implementation reaches only states in which a
\frbl{} node has no live reference.  \texttt{epoch\_run\_safe} is that at the
counter model, so the chain is visible end to end: the counter model refines the
published one, and a client running against it never sees a live reference to a
node it has reclaimed.

The mutation half comes with it, and the way it does is the point.  The ten
heap actions do not involve the implementation at all --- they are the writer's,
they touch the heap and the observation map and never the free list or the
reader set --- so the paired relation admits any step of
\S\ref{sec:inv-mech}'s on the condition that it leaves $R$ and $B$ alone, and
that condition is discharged for all ten by reflexivity: the machine-state
transformers are written to carry $R$ and $B$ through unchanged.  The paired
relation is therefore the union of the two, a run may interleave the writer's
mutations with the protocol freely, and the safety transfer is about whole
programs rather than about the protocol in isolation.

What is \emph{not} claimed is that these are the only steps, or that a
scheduler exists producing any particular interleaving.  The relation says what
may happen, not what does; progress and fairness are the same open questions
they were.

\paragraph{How much of this is about sequential consistency.}
The same question can be asked of the memory model, and it has a sharper
answer.  We are sequentially consistent; Tassarotti et al. verify under
release-acquire, where a thread's view of the heap is a \emph{sub-heap} of what
has been written --- a thread sees some of the writes that have happened and
never one that has not.  So the question a weak reading puts to an invariant is
exactly whether it still holds when the heap shrinks, and an invariant that
survives shrinking transfers from the union of all views, which is what the
writer maintains, to every thread's own view for free.

Asked mechanically: sixteen of the twenty-six conjuncts do not mention the heap
at all and survive any change to it whatever, the proof of each being one
tactic; nine mention it only to rule something out --- the heap occurs in a
hypothesis, or under a negation --- and survive shrinking; and exactly one does
not.

That one is \textbf{HD}.  Every other invariant reads the heap to rule
something out; \textbf{HD} reads it to \emph{require} something, namely that
the node at the end of an edge is in the heap.  A thread that has acquired the
edge but not the node breaks it, and that is not an artificial state: it is the
publication race, the writer initialising a node and then linking it while a
reader acquires the link without the initialisation.  We give it as a witness.

Read as an obligation, \textbf{HD} under a weak memory model \emph{is} the
publish-subscribe guarantee --- Alglave et al.'s second requirement, and the
release-acquire pair Tassarotti et al. name first.  So the single place a
per-thread heap bites is the single place the literature says synchronisation
is required, and it tells us the synchronisation points are not only the
protocol's two but a third at the link, where a writer publishes.  It also says
what is \emph{not} at risk: the memory-safety argument uses \textbf{IFL} and
\textbf{RWOW}, both heap-free, and the publish-subscribe proof uses
\textbf{FNR} and \textbf{FRW}, one heap-free and one that survives shrinking,
so none of the three has anything to say about views.

Stating that as an obligation takes one theorem.  A state is a base together
with a view per thread, each view a sub-heap of the base; \emph{published} says
every thread that can reach a node can see it; and then nineteen of
\textsf{WellFormed}'s twenty conjuncts transfer from the base to every view
with \emph{no hypothesis at all}, the twentieth being exactly that condition.
So a weak account does not re-prove the invariants per thread.  It establishes
one property at the link.

The heap is not the only shared state, and the other kind has the opposite
character.  The registrations are shared too, and it is they that
\S\ref{sec:inv-readside}'s relation ties to the reader set.  For the heap a
thread that sees less is safe; for the registrations a writer that sees
\emph{fewer} than exist concludes a grace period has ended when it has not,
takes the certificate, and frees a node whose snapshot still contains the
reader.  We give that state.  Seeing less is the hazard, not the safeguard.

That is why the two need synchronisation in different places: the heap at the
write that publishes, the registrations at the read that decides the wait is
over.  The relation's first clause is an \emph{iff}, and a sub-map view gives
one direction and loses the one the wait depends on --- so the weak-memory
refinement obligation is not a weakened form of the sequentially consistent one
but the same obligation, with the reads that establish it required to be
acquires.

Those are three synchronisation points, and they are the three Tassarotti et
al.\ name: the link is their Release-Acquire-1, the registration write and the
scan that reads it their second and third.  We arrive at the same three from the
invariants rather than from the algorithm, which is the check on all of this.

The model is written too.  It is release-acquire in its view form: a history
per cell, a view as a timestamp per cell, and three steps --- a write that may
or may not release, \texttt{rcu\_assign\_pointer} being the one that does and
a field initialisation the one that does not, and a read that acquires, which
is \texttt{rcu\_dereference}.  Nothing about views is assumed: the
bookkeeping and the closure release-acquire turns on hold at the initial
configuration and are preserved by every step.

One modelling point is worth naming because the obvious simplification is
wrong.  The released view belongs to the \emph{message}, not to the location.
Attach it to the location and publishing a location twice retroactively changes
what an earlier reader is obliged to have acquired, and the closure is not
preserved.

What comes out is that a reader is at least as far along as the publisher of
anything it can see, and so sees what the publisher wrote --- with the two
programmer's obligations named rather than buried: initialise before you
publish, and do not rewrite afterwards, both of which the \frsh{} discipline
enforces on the write side.  That the release is necessary is not argued but
exhibited: the same run twice, differing in one bit, where the relaxed version
is a perfectly good run of the semantics in which the reader sees the link and
the node's field as it was \emph{before} initialisation.

And then the bridge.  A run under the discipline reaches a configuration in
which every thread's own view satisfies every invariant, given that the
sequentially consistent state whose heap is what has been written does.
Nineteen of the twenty conjuncts come free from the partition above; the
twentieth is publication, and it comes from the semantics.

The registrations are carried by the same three steps, a registration being a
cell, and doing so corrected the guess we started with.  A \emph{stale
generation} read by a reader is safe --- it registers at an earlier generation,
so it looks older than it is and a grace period that need not have waited for
it waits anyway.  A \emph{stale registration} read by the writer is the hazard,
and we give that run.  But reads only move forward, coherence being built into
the read step, so a scan that keeps reading advances monotonically and all it
needs is to reach the last message --- which is the same fairness assumption
the sequentially consistent development already reduces the wait's
\emph{termination} to.  Weak memory therefore adds no new obligation to the
grace period: what makes \textsc{SyncStop} terminate makes it sound.

Reclamation is where the memory model stops being the answer, and that is worth
seeing rather than arguing.  A view holds only what was written, and a thread
that has caught up past the unlink cannot see the old link, so reclamation is
safe exactly when every thread has caught up.  The witness is a reader holding
a reference to a node nothing reachable points at any more, with \emph{both}
writes releasing --- the strongest thing the model has.  No synchronisation
helps; only not reclaiming does.  That is the division of labour the whole
development is about, in one run.

The write-once restriction is lifted, and how that happened is the useful part.
We set out to show it was necessary, by building a reader whose view had never
been the heap.  Every attempt failed at the same step: the acquire pulls the
reader forward \emph{everywhere}, not only at the cell it read.  Under RCU's
actual discipline --- one writer, every write releasing, readers that only read
--- a releasing write publishes the writer's whole knowledge, and with a single
writer those publications are totally ordered, so \emph{every reader's view is
one of the writer's past views}.  A reader is always looking at a heap the
writer really had, and the bridge then needs nothing about which cells are
written twice; it assumes only that the states the writer passed through are
well formed, which is what the rest of this appendix proves.  The restriction
was an artefact of proving the wrong thing.

Two writers and the publications stop being a chain; one relaxed write and the
reader acquires nothing.  Both are outside what \textsc{ToRCUWrite} permits,
which is to say the type system is what enforces the hypothesis of that
theorem.

What stands is that this is release-acquire and not weaker.  Every read
acquires and every view only grows, so there is nothing out of thin air to rule
out, and nothing here would survive a relaxed model unchanged.  A claim about
the Linux kernel memory model, where \texttt{rcu\_dereference} rests on address
dependencies rather than on acquire, would be a different piece of work.

What we claim is narrower than ``RCU is verified under weak memory'' and, we
think, more useful: the single invariant a weak memory model endangers is
publication, publication is what the releasing write buys, and the buying can
be written down.

\paragraph{What is left, and it is three things.}
The third of them used to be ``axiom soundness proper'' and is no longer: all
fifteen actions have it.  What is left in its place is narrower and more
specific, and naming it that way seems more useful than keeping the old
heading.

\emph{Fairness.}  That a reader inside a critical section eventually leaves it.
The wait's termination reduces to it and no further --- it is a property of the
client, and no invariant of the shared state could supply it.

\emph{A fractional heap for the reader's read.}  This is what is left of the
read after the root observation was moved into the structure.  A reader must
learn what a field holds, and the points-to assertion is exclusive.  Nothing
else about the read is open: \texttt{read\_update} discharges the bounding
obligation and grants the observation, and neither it nor the reader's
environment uses a heap fragment at all.  The repair is a fractional or
discardable-fraction points-to, which is standard, and which changes the camera
under the write side as well --- which is why it is named here rather than done
in passing.

\emph{The enumeration premises.}  Axiom soundness itself is done: every action
preserves \textsf{WellFormed}, that is fifteen theorems and now an induction
over runs, and every action also satisfies the denotation of its post-type
environment.  Ten of the fifteen are Hoare triples
--- the whole write side, the binding rules, \textsc{T-Free}, and the two
read-block boundaries, which also compose (\texttt{read\_section\_typed}), so
re-entry is a fact about resources rather than an argument.  The other five are
the step plus the environment with a premise made explicit, and the premises are
of exactly two kinds.

The reader's read has one: the edge it follows, because a reader cannot learn a
cell's contents while the points-to assertion is exclusive.  That is the
fractional heap named above.

The other four --- \textsc{SyncStart}, \textsc{SyncStop}, \textsc{WriteBegin}
and \textsc{WriteEnd} --- have an enumeration premise, and it is the same
premise four times: they are the bulk updates, and a bulk update needs its set
enumerated.  ``Every detached node is stamped'', ``every entry of the writer's
is recoloured'', ``every reachable node is observed'', ``every entry is given
up''.  No resource a thread holds discharges those.  Where it can be stated as a
fact about the thread's \emph{own} map rather than the shared one we state it
that way --- the writer holds its whole column --- which is the honest form and
is what a writer that has tracked its observations since \textsc{WriteBegin}
actually has.  That is what is left, and it is narrower than ``six rules are
missing''.

The read-side denotation is mechanized --- \texttt{D\_rcuItrR}, transcribed from
the figure --- together with the reader's acquisition
(\texttt{reader\_post\_env}).  That one is worth a caution rather than a
heading: the denotation's bounding-thread conjunct asks that a bounding reader's
node already carry a free-list entry, which the reader's own rule does not
establish and which a bounding reader holding an iterator on a live node would
violate.  It appears in the theorem as a hypothesis rather than as something
proved.
